% !TeX root = ../../Book2.tex
% 纲要标题：## 第7章　de Rham、Čech 与层上同调比较
% 纲要位置：几何学计划纲要.md，第 1264 行。
% 正文状态：初稿；尚未正式审核。

% 纲要内容开始
% 
% > 本章目标与证明范围：把三个计算模型及其比较映射真正建立起来。de Rham 部分明确要求光滑流形和形式；本章不使用后置 Hodge 定理来证明 de Rham 后又反向引用。 积分比较所需的单形及单形同伦的相对光滑逼近在本章就地证明，不借后置的一般横截性。
% 
% 纲要内容结束

\chapter{de Rham、Čech 与层上同调比较}
\label{g2:ch:07}
\BookChapterNumber{g2:ch:07}

\begin{chapterabstract}
	局部存在的函数、形式或截面未必能够整体粘合。本章从层的茎和局部正合性出发，以典范消解定义层上同调，证明松弛层、细层的消失及无上同调消解比较。Čech 复形把粘合条件写在开集的交叠上；细化同伦、覆盖比较与覆盖极限说明其与层模型的关系。对光滑流形，就地证明单形及单形同伦的相对光滑逼近，再以积分和双复形建立 de Rham 同构，核对乘法、相对与紧支版本。最后通过周期、射影空间及平坦局部系数，说明实形式记录了哪些拓扑信息、又遗漏了哪些整数信息。
\end{chapterabstract}

\begin{learninggoals}
	\begin{itemize}
		\item 区分预层、层、层化与茎，说明层上的局部正合为什么不保证整体截面正合。
		\item 构造层上同调，证明松弛与细层的消失，并使用无上同调消解计算整体障碍。
		\item 写出 Čech 微分及细化映射，证明指标选择不影响上同调，辨认单份覆盖与覆盖极限的比较条件。
		\item 保持公共面地光滑化奇异单形，证明积分给出自然的 de Rham 同构，并核对乘法和支集。
		\item 计算环面、球面与射影空间的形式类，区分实周期、整周期及挠信息，将局部系统实现为平坦丛的形式复形。
	\end{itemize}
\end{learninggoals}

在圆周的一小段开弧上，总能选择一个连续角度；绕完整个圆周后，角度却可能增加 \(2\pi\)。如果用几段开弧覆盖圆周，每一段的角度都没有问题，问题出现在交叠处：两份角度的差是常数，这些差能否通过修改各片的常数同时消掉？闭一次形式的局部原函数也面临同一个问题。局部求解与整体粘合之间的距离，正是本章要记录的上同调。

\ProseParagraphBreak

第4章把上同调表示为奇异上链，并建立杯积；第5章用局部系统记录绕行运输，第6章用基本类把上同调与同调联系起来。现在我们增加可以直接微分和积分的表示。de Rham 理论\footnote{Georges de Rham（1903--1990），瑞士数学家，研究微分拓扑与代数拓扑，建立了微分形式与拓扑上同调之间的比较。}以闭形式模恰当形式为对象，Čech 理论\footnote{Eduard Čech（1893--1960），捷克数学家，研究射影微分几何与拓扑，发展了以开覆盖组织空间拓扑信息的方法。}以覆盖交叠上的数据为对象，层上同调则允许把同一种局部到整体的问题用于不同类型的系数。对象看起来不同，比较映射必须说明它们为何测量同一障碍。

\begin{table}[htbp]
	\centering
	\caption{本章各模型保存的数据及比较办法}
	\label{g2:tab:cohomology-models}
	\begin{tabular}{lll}
		\toprule
		模型 & 直接保存的数据 & 本章使用的比较 \\
		\midrule
		微分形式 & 局部可微的交替场 & 在光滑单形上积分 \\
		奇异上链 & 每条单形的数值 & 小链与光滑链比较 \\
		Čech 复形 & 各级交叠上的截面 & 细化及覆盖双复形 \\
		层上同调 & 取整体截面后的障碍 & 无上同调消解 \\
		\bottomrule
	\end{tabular}
\end{table}

层态射的满射性给出局部提升，消解把这些提升之间的粘合障碍组织为一个复形；覆盖则将复形写成交叠上的数据。对微分形式，第一卷第14.2节的同伦公式提供局部原形式，第12.3节的单位分解用来组合各片的形式，第14.4节的带角 Stokes 公式使积分与边界相容。第三节先构造保持公共面的光滑逼近，再将这几项局部工具接成积分比较的整体同构。

\ProseParagraphBreak

本章的层论构造先允许一般拓扑空间，覆盖极限比较明确限于 Hausdorff、第二可数、局部紧空间。形式部分取无边界光滑流形和实系数；整数信息需要另由整数上同调比较，不能从微分形式的实向量空间结构中自动恢复。紧支条件也一直单独保留：一个闭紧支形式可能有普通原形式，却没有紧支原形式。

% !TeX root = ../../Book2.tex
% 纲要标题：### 7.1 层与局部正合性
% 纲要位置：几何学计划纲要.md，第 1268 行。
% 正文状态：初稿。

% 纲要内容开始
% 
% * 预层、层化与茎
% * 层的核余核及茎上的正合检测
% * 整体截面与层上同调
% * 松弛层、细层及无上同调机制
% 
% 纲要内容结束

\section{层与局部正合性}
\label{g2:sec:0701}

一个闭一形式在每个小球上都有原函数，原函数却未必能拼成整体函数。此时应同时保存各开集上的解、缩小定义域的限制，以及重叠处的差。层把这些局部资料组织在一起；层上同调则记录局部解无法整体拼合的障碍。本节先把“局部正合”与“整体正合”分开，再构造能够测量两者差别的群。

以下层均为含幺交换环 \(R\) 上的模层；只研究加法时取 \(R=\mathbb Z\)。底空间在需要局部有限覆盖时假设为 Hausdorff、第二可数、局部紧空间，特别包含本书的流形。预层、层化、正合性和松弛层的论证本身不需要这些拓扑条件。

\subsection{预层、层化与茎}
\label{g2:subsec:0701:01}
局部资料先按开集组织：定义域缩小时有一个限制映射，连续缩小两次应与直接限制一致。预层记录这些映射，层再要求相容的局部资料能够唯一粘合。

\begin{definition}\label{g2:def:presheaf-sheaf}
	拓扑空间 \(X\) 上的一个 \(R\) 模\term[yuceng]{预层}[presheaf] \(\mathcal F\)，给每个开集 \(U\) 指定一个模 \(\mathcal F(U)\)，给每个包含 \(V\subseteq U\) 指定限制同态 \(s\mapsto s|_V\)，满足恒等限制及逐次限制律。若对任意开覆盖 \(U=\bigcup_iU_i\)，彼此相容的截面 \(s_i\in\mathcal F(U_i)\) 都能唯一粘成 \(s\in\mathcal F(U)\)，便称 \(\mathcal F\) 为\term[ceng]{层}[sheaf]。空覆盖约定给出 \(\mathcal F(\varnothing)=0\)。与所有限制映射交换的同态族称为预层或层的态射。
\end{definition}

连续函数、光滑函数及光滑微分形式都是层：正则性是局部条件，相同的局部函数又能按其值逐点粘合。相反，“定义域上有界的连续函数”只形成预层；在 \(\mathbb R\) 上用有界开区间覆盖，各处的函数 \(x\mapsto x\) 相容，粘成的函数却不再有界。层公理要求保留任意局部可粘资料，所以有界性这种整体限制通常不应放进层的定义。

\begin{definition}\label{g2:def:stalk}
	对 \(x\in X\)，把所有 \(x\) 的开邻域上的截面作如下识别：\((U,s)\) 与 \((V,t)\) 等价，当且仅当存在 \(x\in W\subseteq U\cap V\)，使 \(s|_W=t|_W\)。所得模
	\[
		\mathcal F_x=\varinjlim_{x\in U}\mathcal F(U)
	\]
	称为 \(\mathcal F\) 在 \(x\) 的\term[jing]{茎}[stalk]，截面 \(s\) 的等价类 \(s_x\) 称为它的芽。加法和数乘在共同邻域上进行，再取等价类。
\end{definition}

芽保留的是某个未指定小邻域内的全部信息，通常远多于点值。例如光滑函数芽 \(f_x\) 决定所有阶导数，\(f(x)=0\) 却不意味着 \(f_x=0\)。后一个等式要求 \(f\) 在某邻域上恒为零。对层而言，两个截面相等当且仅当它们在每点的芽相等：芽相等给出局部相等的覆盖，再用粘合的唯一性。

\begin{proposition}[层化]\label{g2:prop:sheafification}
	任意 \(R\) 模预层 \(\mathcal F\) 有层 \(\mathcal F^{\#}\) 及自然态射 \(\eta:\mathcal F\to\mathcal F^{\#}\)，满足：\(\eta_x\) 在每点都是同构；对任意层 \(\mathcal G\)，每个预层态射 \(\mathcal F\to\mathcal G\) 都唯一经 \(\eta\) 分解。
\end{proposition}
\begin{proof}
	定义 \(\mathcal F^{\#}(U)\) 为芽的族 \((a_x)_{x\in U}\)，要求每个 \(x\in U\) 都有邻域 \(V\subseteq U\) 及 \(s\in\mathcal F(V)\)，使 \(a_y=s_y\) 对所有 \(y\in V\) 成立。限制就是删去定义域外的芽。局部代表可以逐点选择，故相容的芽族唯一粘合；相加时取共同邻域即可，所以它是模层。令 \(\eta(s)=(s_x)_x\)。
	在 \(x\) 附近，一个层化截面按定义来自原预层截面，所以 \(\eta_x\) 满射。若两个原截面的像有同一芽，则在某邻域内它们的芽族一致；尤其在 \(x\) 处原来的芽相等，故 \(\eta_x\) 单射。给定 \(f:\mathcal F\to\mathcal G\)，把局部代表 \(s\) 送至 \(f(s)\)；不同代表在每点有同一芽，因 \(\mathcal G\) 为层，其像在重叠处一致，遂能唯一粘合。局部代表也迫使此延拓唯一。
	因此所得分解由下图表达，其中右下方的箭头由局部代表确定：
	\[
		\begin{tikzcd}[column sep=large]
			\mathcal F \arrow[r,"\eta"] \arrow[dr,"f"'] & \mathcal F^{\#} \arrow[d,dashed,"\exists!\,\bar f"] \\
			& \mathcal G.
		\end{tikzcd}
	\]
\end{proof}

把每个非空开集都赋予同一个模 \(A\)、限制均取恒等，得到常值预层；它的层化记为 \(\underline A\)，称为\term[changzhiceng]{常值层}[constant sheaf]。其截面恰是取值于离散模 \(A\) 的局部常值函数。一个不连通开集上的截面可以在不同分支取不同值，所以“常值层”并不表示每个开集只有一个共同常数。第5章的局部系统也可视为层：在每个平凡化邻域取局部常值截面，以运输给出的过渡映射粘合。

\subsection{层的核余核及茎上的正合检测}
\label{g2:subsec:0701:02}
局部方程 \(f(s)=0\) 可以直接逐开集求解，商却更微妙：一个局部商类可能只能在不同小邻域选不同代表。因此核可以逐开集取，余核则通常要层化。

\begin{proposition}[核、余核与正合性]\label{g2:prop:stalk-exactness}
	设 \(f:\mathcal F\to\mathcal G\) 为拓扑空间 \(X\) 上的 \(R\) 模层态射。核由 \(U\mapsto\ker f_U\) 给出，余核为预层 \(U\mapsto\mathcal G(U)/f\mathcal F(U)\) 的层化；像为预层像的层化。它们在点 \(x\) 的茎分别是 \(\ker f_x\)、\(\operatorname{coker}f_x\) 和 \(\operatorname{im}f_x\)。因而层的复形正合，当且仅当每点的茎复形正合；层态射为同构，当且仅当所有茎映射为同构。
\end{proposition}
\begin{proof}
	核截面的相容族在 \(\mathcal F\) 中粘合，像在 \(\mathcal G\) 中局部为零，故整体也为零；这证明核已是层。商的泛性质先在预层中成立，再用\cref{g2:prop:sheafification}得到层范畴中的余核泛性质。
	关键是取芽保持正合。若 \(s_x\) 在某映射下为零，就能缩小代表 \(s\) 的定义域，使其像作为截面已经为零。若一个芽属于某像，选择该原像芽的代表，再缩小邻域便得到截面层面的原像。有限次等式和有限个代表只需取有限个邻域的交；这证明有向极限保持模的核、像及余核。层化不改茎，故上述茎公式成立。
	若每点核与像相同，则核截面在每点附近来自原像，正是像层的定义；反向可直接取茎。最后，茎同构意味着核层与余核层的所有茎为零，而所有茎为零的层只能是零层；也可以把局部逆像利用唯一性直接粘成逆态射。
\end{proof}

特别地，满的层态射只保证截面\emph{局部}有原像，不保证每个开集上的截面映射满射。后文写
\[
	0\longrightarrow\mathcal F\longrightarrow\mathcal G\longrightarrow\mathcal H\longrightarrow0
\]
时，正合性均按\cref{g2:prop:stalk-exactness}理解。层化使商对象补足局部代表，因而核与余核满足通常的模运算规则；涉及核、像的图追逐可以逐茎进行，但追逐所得的局部原像仍可能没有整体选择。

\subsection{整体截面与层上同调}
\label{g2:subsec:0701:03}
记 \(\Gamma(U,\mathcal F)=\mathcal F(U)\)，特别称 \(\Gamma(X,\mathcal F)\) 为整体截面模。对上述短正合列，
\[
	0\longrightarrow\Gamma(X,\mathcal F)\longrightarrow\Gamma(X,\mathcal G)\longrightarrow\Gamma(X,\mathcal H)
\]
总是正合：映到零的整体截面在局部属于核，核本身已经是层，故各局部原像唯一粘合。这叫整体截面的\term[zuozhenghe]{左正合性}[left exactness]，最后一箭却可能不满。

\begin{example}\label{g2:exm:sheaf-antiderivative-obstruction}
	在 \(S^1\) 上，令 \(\mathcal C^{\infty}\) 为实光滑函数层，\(\Omega^1\) 为实一形式层。局部区间上的积分说明
	\[
		0\longrightarrow\underline{\mathbb R}\longrightarrow\mathcal C^{\infty}\xrightarrow{\mathrm d}\Omega^1\longrightarrow0
	\]
	在茎上正合。可是若把 \(S^1\) 写成 \(\mathbb R/\mathbb Z\)，形式 \(\mathrm dt\) 沿一周的积分为 \(1\)，而任何 \(\mathrm df\) 沿一周的积分为 \(0\)。所以整体截面上的最后一箭不满。局部原函数在重叠处只相差常数；这些常数无法全被重新选择的局部常数消去，才是整体障碍。
\end{example}

为记录局部提升的粘合障碍，先允许在每点独立指定一个芽，再取商记录这些任意芽与真实局部截面之间的差别。反复进行这两项操作，就能构造一条具有自然映射的消解。对层 \(\mathcal F\)，设
\[
	T\mathcal F(U)=\prod_{x\in U}\mathcal F_x.
\]
这里不要求芽随点连续，也不要求它们由同一个局部截面产生；限制只是投影。故 \(T\mathcal F\) 是层，且 \(\mathcal F\to T\mathcal F\)，\(s\mapsto(s_x)_x\)，为单射。取层余核 \(Q\mathcal F=T\mathcal F/\mathcal F\)，再迭代，令
\begin{equation}\label{g2:eq:canonical-sheaf-resolution}
	G^q\mathcal F=TQ^q\mathcal F,\qquad
	0\longrightarrow\mathcal F\longrightarrow G^0\mathcal F\xrightarrow{\mathrm d_G}G^1\mathcal F\xrightarrow{\mathrm d_G}\cdots,
\end{equation}
其中 \(Q^0\mathcal F=\mathcal F\)，微分先取商 \(TQ^q\mathcal F\to Q^{q+1}\mathcal F\)，再嵌入 \(TQ^{q+1}\mathcal F\)。像与核都等于中间的 \(Q^{q+1}\mathcal F\)，所以这是一条正合的层复形，称为本节的典范消解。

\begin{definition}\label{g2:def:sheaf-cohomology}
	拓扑空间 \(X\) 上的 \(R\) 模层 \(\mathcal F\) 的\term[cengshangtongdiao]{层上同调}[sheaf cohomology]定义为
	\[
		H^q(X,\mathcal F)=H^q\bigl(\Gamma(X,G^\bullet\mathcal F)\bigr),\qquad q\ge0.
	\]
	若对所有 \(q>0\) 都有 \(H^q(X,\mathcal F)=0\)，称 \(\mathcal F\) 在 \(X\) 上\term[wushangtongdiao]{无上同调}[acyclic]。
\end{definition}

零次核恰是 \(\Gamma(X,\mathcal F)\)，故 \(H^0(X,\mathcal F)=\Gamma(X,\mathcal F)\)。在开集 \(U\) 上限制时，茎、层商及上述逐点乘积都与限制交换，因此 \(G^\bullet\mathcal F|_U=G^\bullet(\mathcal F|_U)\)。这既定义了上同调的限制，也说明 \(H^q(U,\mathcal F)\) 无歧义。

\begin{proposition}[层上同调的长正合列]\label{g2:prop:sheaf-long-exact}
	拓扑空间 \(X\) 上每条 \(R\) 模层短正合列 \(0\to\mathcal F\to\mathcal G\to\mathcal H\to0\) 都给出自然长正合列
	\[
		0\to H^0(X,\mathcal F)\to H^0(X,\mathcal G)\to H^0(X,\mathcal H)
		\xrightarrow{\partial}H^1(X,\mathcal F)\to\cdots.
	\]
	连接同态 \(\partial s\) 为零当且仅当 \(s\in\Gamma(X,\mathcal H)\) 能整体提升到 \(\mathcal G\)。
\end{proposition}
\begin{proof}
	茎序列正合，而模的任意乘积保持短正合：满射的每个坐标各选一个原像即可。因此 \(T\) 以及 \(\Gamma(X,T(-))\) 都保持短正合。把短正合列与其 \(T\) 像排成两行，竖箭为单射；逐茎追逐得到余核列 \(0\to Q\mathcal F\to Q\mathcal G\to Q\mathcal H\to0\) 正合。更具体地，右侧一个商类先提升到中间 \(T\) 项，差额来自左侧 \(T\) 项；若中间商类落到零，再减去中间原层的局部提升，便回到左侧商类。这也证明 \(Q\) 正合，迭代后 \(Q^q\) 正合。
	故 \(\Gamma(X,G^\bullet(-))\) 将原序列变成逐次短正合的上链复形序列。连接映射的构造与第2章的链版本相同：将右侧闭上链提升至中间，对其取微分，再由左侧单射识别所得闭上链。两种提升相差中间的左侧上链，其微分只改变一个边界。若连接类为零，减去该左侧原像后，中间提升便是闭上链；反之闭提升的微分为零。这证明连接处的正合性，另外两处由同样的提升和取像直接得到。所有步骤与短正合列的态射交换，故自然。零次的结论正是最后一句。
\end{proof}

每个层态射都逐点作用于芽，并诱导层商上的态射，因而自动给出典范消解之间的链映射。接下来先考察使高次上同调消失的延拓和局部组合机制；下一节再据此用更适合计算的消解替换典范消解。

\subsection{松弛层、细层及无上同调机制}
\label{g2:subsec:0701:04}
\begin{definition}\label{g2:def:flabby-fine}
	若层 \(\mathcal F\) 对所有开集包含 \(V\subseteq U\) 的限制 \(\mathcal F(U)\to\mathcal F(V)\) 都满射，称它为\term[songchiceng]{松弛层}[flabby sheaf]。在仿紧 Hausdorff 空间 \(X\) 上，若对每个局部有限开覆盖 \(\{U_i\}\)，都有层端同态 \(\rho_i:\mathcal F\to\mathcal F\) 及局部有限闭集族 \(K_i\subseteq U_i\)，使 \(\rho_i|_{X\setminus K_i}=0\) 且 \(\sum_i\rho_i=1\)，则称 \(\mathcal F\) 为\term[xiceng]{细层}[fine sheaf]。端同态之和在每个足够小的邻域都是有限和。
\end{definition}

松弛性允许任意延拓；细性允许把局部资料乘上有支集控制的权重后相加。两者解决障碍的方式不同。\(T\mathcal F\) 总是松弛，因为一族芽可以在较大开集其余各点取零。光滑函数层通常不松弛，例如 \(1/x\) 在 \((0,1)\) 上无法延拓成 \(\mathbb R\) 上的光滑函数；但它是细层。

\begin{lemma}[松弛核的提升性质]\label{g2:lem:flabby-lifting}
	若拓扑空间 \(X\) 上的层短正合列 \(0\to\mathcal A\to\mathcal B\to\mathcal C\to0\) 的核 \(\mathcal A\) 松弛，则 \(\mathcal B(U)\to\mathcal C(U)\) 对每个开集 \(U\) 满射。若 \(\mathcal B\) 也松弛，则 \(\mathcal C\) 松弛。
\end{lemma}
\begin{proof}
	固定 \(s\in\mathcal C(U)\)，考虑它在开子集 \(V\subseteq U\) 上的提升 \(b_V\)，按延拓排序。全序族可由层公理在开集之并上粘合，故 Zorn 引理给出极大提升。若 \(x\in U\setminus V\)，茎上的满射性给出邻域 \(W\) 上的提升 \(b_W\)。在 \(V\cap W\) 上，两者之差来自 \(\mathcal A(V\cap W)\) 的唯一截面；利用松弛性把这个差延拓至 \(W\)，从 \(b_W\) 中减去后便与 \(b_V\) 相合。因此提升能延至 \(V\cup W\)，违背极大性。所以 \(V=U\)。若 \(\mathcal B\) 也松弛，先把 \(\mathcal C(V)\) 的截面提升到 \(\mathcal B(V)\)，再延至 \(\mathcal B(U)\) 并取像，就得到 \(\mathcal C\) 的松弛性。
\end{proof}

\begin{theorem}[松弛层与细层的消失]\label{g2:thm:flabby-fine-acyclic}
	拓扑空间 \(X\) 上的松弛 \(R\) 模层无上同调，其在每个开子集上的限制也无上同调。仿紧 Hausdorff 空间 \(X\) 上的细 \(R\) 模层在 \(X\) 上无上同调。
\end{theorem}
\begin{proof}
	若 \(\mathcal F\) 松弛，则在 \(0\to\mathcal F\to T\mathcal F\to Q\mathcal F\to0\) 中，前两项都松弛。\cref{g2:lem:flabby-lifting}表明 \(Q\mathcal F\) 松弛，且取整体截面仍正合。逐次迭代，典范消解的整体截面复形在正次数均正合，故上同调消失。松弛性显然保留于开子集上的限制，同一论证遂适用。
	现设 \(\mathcal F\) 细，取 \(q>0\) 及 \(z\in\Gamma(X,G^q\mathcal F)\) 满足 \(\mathrm d_Gz=0\)。层消解的正合性给出一个开覆盖，使每个开集上存在 \(b_i\) 满足 \(\mathrm d_Gb_i=z\)。用仿紧性取局部有限开细化，仍记为 \(\{U_i\}\)，并选细性给出的 \(\rho_i,K_i\)。函子性使 \(\rho_i\) 诱导消解端同态 \(G^\bullet\rho_i\)，与 \(\mathrm d_G\) 交换。由于消解构造与开限制交换，\(G^\bullet\rho_i\) 在 \(X\setminus K_i\) 上也为零。
	因此 \((G^{q-1}\rho_i)b_i\) 可在 \(U_i\) 外补零：在 \(U_i\cap(X\setminus K_i)\) 上两个定义都是零，层公理保证粘合。令 \(b\) 为这些整体截面的局部有限和。每点附近只有有限项，故加法、函子作用及微分均与此和交换，得到
	\[
		\mathrm d_Gb=\sum_i(G^q\rho_i)z
		=G^q\Bigl(\sum_i\rho_i\Bigr)z=z.
	\]
	每个正次闭上链都成为边界，证明细层无上同调。这里没有要求一般微分与任意函数相乘交换；交换的是同一个层端同态在\emph{典范消解}上诱导的映射。
\end{proof}

\begin{corollary}\label{g2:cor:smooth-module-acyclic}
	有限维、Hausdorff、第二可数的光滑流形 \(M\) 上，每个光滑函数层 \(\mathcal C^{\infty}_M\) 的模层都是细层，因而无上同调。特别地，光滑向量丛截面层及微分形式层 \(\Omega^p_M\) 均无上同调；它们在任意开子集上的限制也无上同调。
\end{corollary}
\begin{proof}
	第一卷定理12.3.5说明：对每个开覆盖，可取光滑函数 \(\chi_i\)，支集族局部有限、\(\operatorname{supp}\chi_i\subseteq U_i\)，且 \(\sum_i\chi_i=1\)。在模层上取 \(\rho_i(s)=\chi_i s\)，便满足\cref{g2:def:flabby-fine}。流形的每个开子集仍为 Hausdorff、第二可数光滑流形，故限制后的结论同样成立。
\end{proof}

应用于\cref{g2:exm:sheaf-antiderivative-obstruction}，长正合列与 \(\mathcal C^{\infty}\) 的消失给出
\[
	H^1(S^1,\underline{\mathbb R})
	\cong\Omega^1(S^1)/\mathrm d\mathcal C^{\infty}(S^1)\cong\mathbb R.
\]
最后一个同构由沿一周积分给出。若 \(a(t)\mathrm dt\) 的积分为零，\(f(t)=\int_0^t a(u)\,\mathrm du\) 是周期函数，故在圆上给出原函数；这证明其核恰为恰当形式。于是层上同调确实把局部原函数的粘合失败转成了可以计算的整体不变量。
% 纲要条目排为编号小节标题，并进入目录与书签；说明段仍保留于纲要注释中。

% !TeX root = ../../Book2.tex
% 纲要标题：### 7.2 消解和 Čech 模型
% 纲要位置：几何学计划纲要.md，第 1225 行。
% 正文状态：未撰写；本文件目前是框架。

% 纲要内容开始
% 
% * 无上同调消解及比较定理
% * 覆盖的 Čech 复形
% * 细化映射与覆盖选择
% * 无上同调覆盖的比较定理
% 
% 纲要内容结束

\section{消解和 Čech 模型}
\label{g2:sec:0702}

待撰写

\subsection{无上同调消解及比较定理}
\label{g2:subsec:0702:01}
待撰写

\subsection{覆盖的 Čech 复形}
\label{g2:subsec:0702:02}
待撰写

\subsection{细化映射与覆盖选择}
\label{g2:subsec:0702:03}
待撰写

\subsection{无上同调覆盖的比较定理}
\label{g2:subsec:0702:04}
待撰写
% 纲要条目排为编号小节标题，并进入目录与书签；说明段仍保留于纲要注释中。

% !TeX root = ../../Book2.tex
% 纲要标题：### 7.3 积分与 de Rham 定理
% 纲要位置：几何学计划纲要.md，第 1282 行。
% 正文状态：初稿；尚未正式审核。

% 纲要内容开始
% 
% * 光滑奇异链及光滑逼近
% * 积分映射与 Stokes
% * 好覆盖和局部 Poincaré 引理
% * 双复形的两种收缩
% * 乘法和支集版本的比较
% 
% 纲要内容结束

\section{积分与 de Rham 定理}
\label{g2:sec:0703}

本节把局部形式与拓扑上链联系起来。第一卷已经构造外微分、形式积分及 Stokes 公式；这里先证明连续单形能够在保留公共面的条件下光滑化，再用覆盖双复形证明积分诱导同构。整个比较均在有限维、Hausdorff、第二可数、无边界光滑流形上进行。

\subsection{光滑奇异链及光滑逼近}
\label{g2:subsec:0703:01}
奇异单形只要求连续，而拉回微分形式需要导数。因此，积分以前先要解决一个几何问题：能否把用于表示同调类的有限个单形同时光滑化，并使公共面仍然完全相同？只把每个单形单独逼近，原先抵消的边界可能不再抵消；保留公共面是下面构造的核心。

\begin{definition}[光滑奇异链]\label{g2:def:smooth-singular-chain}
	设 \(M\) 为有限维、Hausdorff、第二可数的无边界光滑流形，\(R\) 为含幺交换环。映射 \(\sigma:\Delta^k\to M\) 若可延拓为标准闭单形在其仿射包中某个开邻域上的光滑映射，就称为\term[guanghua qiyi danxing]{光滑奇异单形}[smooth singular simplex]。以这些映射为基的自由 \(R\)-模记为 \(C_k^\infty(M;R)\)。面映射 \(\delta_i:\Delta^{k-1}\to\Delta^k\) 是仿射映射，故
	\[
		\partial\sigma=\sum_{i=0}^k(-1)^i\sigma\circ\delta_i
	\]
	仍是光滑链。所得同调记为 \(H_k^\infty(M;R)\)，光滑奇异上链复形记为 \(C_\infty^k(M;R)=\operatorname{Hom}_R(C_k^\infty(M;R),R)\)。
\end{definition}

这里要求在闭单形附近延拓，因而边、顶点处也有足够的导数。单形本身有角，并不要求它的像是嵌入子流形；常值单形和微分降秩的单形都允许。光滑链的构造与比较可参阅 \cite[第18章，Smooth Singular Homology]{Lee2012Smooth}。下面把所需的相对逼近直接证明出来。

\begin{lemma}[紧像附近的欧氏模型]\label{g2:lem:compact-image-retraction}
	设 \(M\) 为有限维、Hausdorff、第二可数的无边界光滑流形，\(Q\subset M\) 紧。存在 \(Q\) 的开邻域 \(V\)、光滑嵌入 \(e:V\to\mathbb R^N\)、包含 \(e(Q)\) 的开集 \(W\subset\mathbb R^N\) 和光滑映射 \(r:W\to V\)，使 \(r(e(x))=x\) 对 \(Q\) 的某个开邻域中的 \(x\) 成立。
\end{lemma}
\begin{proof}
	第一卷推论12.3.6的闭集截断结论允许在相对紧坐标图 \((U_i,\varphi_i)\) 内选紧支光滑函数 \(\chi_i\)，使有限个集合 \(\{\chi_i>0\}\) 覆盖 \(Q\)。令 \(V=\{\sum_i\chi_i>0\}\)，定义
	\[
		e(x)=\bigl(\chi_1(x),\chi_1(x)\varphi_1(x),\ldots,\chi_s(x),\chi_s(x)\varphi_s(x)\bigr).
	\]
	每项在相应坐标图之外补零；支集紧含于图内保证补零光滑。如果 \(\chi_i(x)>0\)，从第 \(i\) 组分量的商便恢复 \(\varphi_i(x)\)。这同时证明 \(e\) 单射、逆映射在像上连续，以及 \(\mathrm{d}e_x\) 单射，所以 \(e\) 是嵌入。

	说明为什么在 \(e(Q)\) 附近能回到这张嵌入像上。在每个 \(x\in V\) 处取 \(\mathrm{d}e_x(\mathrm{T}_xV)\) 的欧氏正交补，考虑
	\[
		\Psi(x,v)=e(x)+v,\qquad v\perp\mathrm{d}e_x(\mathrm{T}_xV).
	\]
	这些正交补组成光滑向量丛：在坐标图内，切空间由一个满秩矩阵 \(A(x)\) 的列生成，其正交投影为 \(A(A^{\mathsf T}A)^{-1}A^{\mathsf T}\)，随 \(x\) 光滑变化。在零截面处，\(\mathrm{d}\Psi\) 把切向与法向分量相加，是同构。第一卷定理11.3.3的逆函数定理因而给出每个 \((x,0)\) 附近的光滑逆。

	选 \(Q\subset V_0\Subset V\)。在基点属于 \(V_0\)、\(|v|\) 充分小时，\(\Psi\) 可以同时取为单射。否则存在两组不同的 \((x_j,v_j),(y_j,w_j)\)，基点属于 \(\overline V_0\)，法向长度趋于零，而 \(\Psi(x_j,v_j)=\Psi(y_j,w_j)\)。取收敛子列后有 \(e(x)=e(y)\)，故 \(x=y\)；两组点最终都落入 \((x,0)\) 附近同一逆函数邻域，与其单射性矛盾。类似地，紧性使微分在足够小的这一邻域处处可逆。因此 \(\Psi\) 从这个开邻域到其开像是微分同胚，投影其逆映射的第一分量便得到 \(r\)。
\end{proof}

平滑目标为向量空间时，面间兼容可以用一个直接的延拓公式处理。这个步骤也解释了为什么“每个面上光滑”在单形的角处足够。

\begin{lemma}[单形边界数据的光滑延拓]\label{g2:lem:simplex-face-extension}
	设 \(b:\partial\Delta^k\to\mathbb R^N\) 连续，且在每个闭面上都可光滑延拓。则存在 \(\Delta^k\) 的开邻域上的光滑映射 \(B\)，使 \(B|_{\partial\Delta^k}=b\)。若连续映射 \(F:\Delta^k\to\mathbb R^N\) 满足 \(F|_{\partial\Delta^k}=b\)，则对每个 \(\varepsilon>0\)，可取这样的 \(B\) 使 \(\sup_{\Delta^k}|B-F|<\varepsilon\)。
\end{lemma}
\begin{proof}
	先只做边界延拓。在一个边界点附近，把经过该点的 \(q\) 个面写成 \(x_1=0,\ldots,x_q=0\)，单形位于 \(x_i\geq0\) 一侧；未经过该点的面可通过缩小邻域避开。对非空 \(I\subseteq\{1,\ldots,q\}\)，记 \(p_I\) 为把这些坐标置零的投影，\(b_I\) 为 \(b\) 在相应交面上的限制，并选其光滑延拓 \(\widetilde b_I\)。在足够小的邻域内定义
	\begin{equation}\label{g2:eq:corner-face-extension}
		B_0(x)=\sum_{\varnothing\ne I\subseteq\{1,\ldots,q\}}(-1)^{|I|+1}\widetilde b_I(p_Ix).
	\end{equation}
	若 \(x\) 位于第 \(j\) 个面内，含 \(j\) 与不含 \(j\) 的项两两抵消，只留下 \(b(x)\)。这里所有投影点仍在相应闭面内，所以不同延拓在面外怎样选择并不影响这次抵消。用边界附近的有限光滑单位分解粘合这些 \(B_0\)，再乘一个在边界附近等于一的截断函数并补零，得到整个单形邻域上的光滑 \(B_1\)。

	现在令 \(h=F-B_1\)。它连续且在边界为零。由一致连续性，\(h\) 在足够窄的边界邻域内任意小。取光滑截断 \(\chi\)，支集紧含于单形内部，在这个窄邻域之外等于一，则 \(\chi h\) 在全仿射包上补零仍连续，并满足 \(\|\chi h-h\|_\infty<\varepsilon/2\)。选非负光滑核 \(\rho_\delta\)，支集在半径 \(\delta\) 的球内且积分为一。卷积 \(h_\delta=\rho_\delta*(\chi h)\) 光滑；当 \(\delta\) 足够小时，它的支集仍紧含于单形内部，而且
	\[
		\|h_\delta-\chi h\|_\infty
		\leq\sup_{|y|\leq\delta,\,x}|\chi(x-y)h(x-y)-\chi(x)h(x)|<\varepsilon/2.
	\]
	因此 \(B=B_1+h_\delta\) 满足要求。零维时没有边界，直接取给定点值即可。
\end{proof}

\begin{lemma}[相对单形光滑逼近]\label{g2:lem:relative-simplex-smoothing}
	设 \(K\) 是有限单形复形，\(A\subset K\) 为子复形，\(M\) 为有限维、Hausdorff、第二可数的无边界光滑流形。若 \(f:|K|\to M\) 连续，且 \(f\) 在 \(A\) 的每个闭单形上光滑，则存在逐单形光滑的 \(g:|K|\to M\) 和相对 \(|A|\) 的同伦 \(H:f\simeq g\)。可以要求 \(H\) 任意接近 \(f\)：对每个单形预先指定包含其像的开集，整个单形同伦仍留在该开集内。

	结论也适用于有限个单形按面识别得到的复形，包括同一个单形的不同面之间的识别。若 \(f\) 在某个单形上原已光滑，则 \(H\) 在该单形与 \([0,1]\) 的乘积上可取光滑。给定一条连续单形同伦，若其两端及指定侧面已逐面光滑，则可保持它们不变，将同伦光滑化于棱柱的有限单形剖分上。
\end{lemma}
\begin{proof}
	把紧像 \(f(|K|)\) 放入\cref{g2:lem:compact-image-retraction}的模型中，记 \(F=e\circ f\)。先在 \(A\) 上保持 \(G=F\)，再依维数把 \(G\) 定义到其余单形。假设一个 \(k\) 维单形的边界上已经给定逐面光滑的 \(b\)，而且 \(b\) 接近 \(F\)。边界误差 \(b-F\) 可连续延拓到整个单形而不增加上界：从重心 \(a\) 出发，把 \(x\ne a\) 唯一写成 \(x=(1-t)a+ty\)，其中 \(y\in\partial\Delta^k\)、\(0<t\leq1\)，令
	\[
		E(x)=t\bigl(b(y)-F(y)\bigr),\qquad E(a)=0.
	\]
	这个延拓连续，\(F+E\) 恰好限制为 \(b\)。应用\cref{g2:lem:simplex-face-extension}，得到一个光滑 \(G\)，在边界上严格等于 \(b\)，并可任意接近 \(F+E\)。若第 \(k\) 维步骤另添的误差小于 \(\varepsilon/(\dim K+1)\)，全部步骤结束后的误差小于 \(\varepsilon\)。同维单形内部互不影响，公共面在处理它们以前已经固定，因此这些映射粘合为连续 \(G\)。对于按面识别的有限单形族，给每个已识别的面只选择一次数据，同一构造仍成立。

	取 \(\varepsilon\) 充分小，使连接 \(F(x)\) 与 \(G(x)\) 的所有线段均落在回缩的定义域 \(W\) 内。定义
	\begin{equation}\label{g2:eq:relative-smoothing-homotopy}
		g(x)=r(G(x)),\qquad H(x,t)=r\bigl((1-t)F(x)+tG(x)\bigr).
	\end{equation}
	在 \(A\) 上两端完全相等，故同伦固定 \(A\)。紧性以及 \(r\) 的连续性保证：进一步缩小误差后，每个单形的整条同伦都留在其指定开集内。若 \(f\) 在一个单形上原已光滑，公式中的 \(F\)、\(G\) 都有邻域延拓，因而 \(H\) 在该单形的棱柱上光滑。

	最后，把 \(\Delta^k\times[0,1]\) 按第2章的棱柱剖分写成有限个单形。若一条连续同伦的两端和指定侧面已经逐面光滑，就将这些面作为子复形 \(A\)，再次应用刚证的构造。得到的同伦与原同伦有相同的两端及指定侧面，并在每个棱柱单形上光滑。它在剖分界面上不必具有来自两个单形的同一法向导数；积分链只要求各闭单形光滑且公共面相同，这正是所需的相容条件。
\end{proof}

\begin{theorem}[光滑链与连续链的比较]\label{g2:thm:smooth-singular-comparison}
	设 \(M\) 为有限维、Hausdorff、第二可数的无边界光滑流形，\(R\) 为含幺交换环。包含映射诱导自然同构
	\[
		H_k^\infty(M;R)\xrightarrow{\ \cong\ }H_k(M;R).
	\]
	特别地，限制实上链得到自然同构
	\[
		H^k(M;\mathbb R)\xrightarrow{\ \cong\ }H^k(C_\infty^\bullet(M;\mathbb R)).
	\]
	这里的自然性指光滑映射之间的自然性。
\end{theorem}
\begin{proof}
	每条链只涉及有限个奇异单形。把它们以及全部面列出，相同的参数映射只保留一次，并按面关系识别其定义域。这给出上引理允许的有限单形族；即使两个不同面恰好是同一个奇异单形，它们也得到同一个光滑替代和同伦。以 \(S\) 表示这个有限族上的光滑替代，面相容性给出 \(\partial S=S\partial\)。

	设 \(z\) 是连续奇异圈。对其有限单形族应用上引理，则 \(Sz\) 为光滑圈。用同伦棱柱定义 \(P\)，第2章\cref{g2:eq:prism-identity}的边界计算逐个单形给出
	\[
		\partial P+P\partial=S-\operatorname{id}.
	\]
	原公式只使用相邻面上的同伦相同，故适用于这里逐单形给出的相容同伦。于是 \(Sz-z=\partial Pz\)，证明满射。若光滑圈 \(z\) 在连续链中满足 \(z=\partial b\)，把 \(b\)、\(z\) 及其全部面同时列出，以 \(z\) 的单形及其面为固定子复形。相对逼近保持这些单形不变，故 \(\partial Sb=Sz=z\)，证明单射。这个论证使用的都是有限 \(R\)-线性组合，不要求链的系数为整数。

	为说明实上同调的结论，回顾向量空间复形的直接事实：限制到圈诱导
	\[
		H^k(\operatorname{Hom}_{\mathbb R}(C_*,\mathbb R))
		\cong\operatorname{Hom}_{\mathbb R}(H_k(C_*),\mathbb R).
	\]
	每个同调上的线性泛函先给出圈上的泛函，再经选取基延拓到 \(C_k\)，因其消去边界而成为余圈；反过来，一个余圈若在所有圈上为零，就经 \(\partial:C_k\to B_{k-1}\) 因子化，把该泛函从 \(B_{k-1}\) 延拓到 \(C_{k-1}\) 后，它便是余边界。这证明该识别，且求值公式说明识别自然。将已证同调同构代入，便得实上同调同构。光滑映射把光滑单形送到光滑单形，与包含及求值交换，故自然性成立。
\end{proof}

\begin{corollary}[光滑小链与开集对]\label{g2:cor:smooth-small-chains}
	设 \(M\) 为有限维、Hausdorff、第二可数的无边界光滑流形，\(\mathcal U\) 为其开覆盖。由像落在某个 \(U\in\mathcal U\) 内的光滑单形生成的子复形 \(C_*^{\infty,\mathcal U}(M;R)\) 包含于 \(C_*^\infty(M;R)\)，并诱导同调同构。若 \(U\subset M\) 开，则光滑相对链与连续相对链也计算同一个 \(H_*(M,U;R)\)；实系数相对上同调亦然。
\end{corollary}
\begin{proof}
	\cref{g2:prop:subdivision-chain-map,g2:prop:subdivision-chain-homotopy}中的细分及链同伦，只在参数单形内使用仿射映射与仿射锥。因此，原单形光滑时，产生的每个单形仍光滑，而且像始终留在原单形的像中。有限光滑链经足够多次细分成为小链，细分链同伦给出满射；若小圈在全部光滑链中有填充，把这个填充细分为小链，再利用原小圈与其细分之间仍属小链的链同伦，便得到小链填充。这正是\cref{g2:thm:small-chain-isomorphism}的证明，其中所有链现在都保留光滑性。

	开子集 \(U\) 也是无边界光滑流形。比较 \((M,U)\) 的两条相对同调长正合列，绝对项由上一定理同构，五引理给出相对同构；实上同调再用证明中的线性泛函识别。这里要求 \(U\) 开，是为了让光滑单形在 \(U\) 中的邻域延拓与在 \(M\) 中的定义相符。
\end{proof}

\subsection{积分映射与 Stokes}
\label{g2:subsec:0703:02}
以下 \(M\) 为有限维、Hausdorff、第二可数、无边界的光滑流形，形式取实系数。写 \(C_\infty^k(M;\mathbb R)=\operatorname{Hom}_{\mathbb R}(C_k^\infty(M;\mathbb R),\mathbb R)\)，上边界仍为 \((\delta\varphi)(c)=\varphi(\partial c)\)。一个形式给上链赋值的办法已经很自然：把它在每条光滑单形上积分。

\begin{definition}[积分上链]\label{g2:def:integration-cochain}
	对 \(\omega\in\Omega^k(M)\)，定义
	\[
		I_M\omega(\sigma)=\int_{\Delta^k}\sigma^*\omega.
	\]
	标准单形采用有序顶点 \((0,e_1,\ldots,e_k)\) 及定向 \(\mathrm{d}t^1\wedge\cdots\wedge\mathrm{d}t^k\)；零次积分是函数在点上的值。此定义按线性延拓到链上。
\end{definition}

这里不要求 \(\omega\) 紧支，因为每条积分单形的定义域紧。第一卷定理14.4.5的带角 Stokes 公式\footnote{George Gabriel Stokes（1819--1903），生于爱尔兰的英国数学家、物理学家，研究流体力学与光学。}适用于闭单形：每个面带上边界定向，恰好产生奇异边界中的交替符号。因此对 \((k+1)\) 单形 \(\sigma\)，
\begin{equation}\label{g2:eq:integration-cochain-map}
	I_M(\mathrm{d}\omega)(\sigma)
	=\int_{\Delta^{k+1}}\mathrm{d}(\sigma^*\omega)
	=\sum_{j=0}^{k+1}(-1)^j\int_{\Delta^k}(\sigma\circ d_j)^*\omega
	=(\delta I_M\omega)(\sigma).
\end{equation}
第一卷定理14.5.5已证明这一链级等式；现在\cref{g2:thm:smooth-singular-comparison}使它能够取值于通常的奇异上同调。具体说，光滑链包含诱导上链限制 \(j^*:C^*(M;\mathbb R)\to C_\infty^*(M;\mathbb R)\)，它在上同调上为同构，故积分定义
\[
	\mathcal I_M:H_{\mathrm{dR}}^k(M)\longrightarrow H^k(M;\mathbb R),
	\qquad \mathcal I_M=(H^k(j^*))^{-1}\circ H^k(I_M).
\]
\symidx{\(\mathcal I_M\)：由光滑单形积分定义的 de Rham 比较映射}
上链限制在实数域上保持同调同构，可直接用万有系数定理的域版本说明：向量空间复形可选取边界、同调及其补空间，取线性对偶后仅留下同调的对偶；不需要这些空间有限维。

\begin{proposition}[积分的自然性]\label{g2:prop:integration-naturality}
	若 \(f:M\to N\) 光滑，则 \(I_Mf^*=f^*I_N\)。因此 \(\mathcal I\) 与光滑拉回相容，特别与开子集限制相容。若 \(z\) 是光滑 \(k\) 圈，则
	\[
		\langle\mathcal I_M[\omega],[z]\rangle=\int_z\omega.
	\]
\end{proposition}
\begin{proof}
	在任意光滑单形 \(\sigma\) 上，\(\sigma^*f^*\omega=(f\circ\sigma)^*\omega\)，所以两个上链的值相等。光滑链包含也与 \(f\) 相容，故在其诱导同构两侧取逆仍得到自然性。求值公式正是定义。若改变闭形式代表或圈代表，差分别由\eqref{g2:eq:integration-cochain-map}化为对圈的余边界求值或对边界的闭上链求值，均为零。
\end{proof}

例如圆周上写 \(\theta=(x\,\mathrm{d}y-y\,\mathrm{d}x)/(2\pi)\)，其中 \(x^2+y^2=1\)。沿 \(t\mapsto(\cos2\pi t,\sin2\pi t)\) 的积分为 \(1\)，所以 \(\mathcal I[\theta]\) 检测一次正向绕行。这不仅说明 \(\theta\) 闭而非恰当，还固定了它与整数生成元比较时的归一化。

\subsection{好覆盖和局部 Poincaré 引理}
\label{g2:subsec:0703:03}
局部求原形式与整体求原形式的差别，可以在最简单的零次就看见：每个连通小开集上的闭函数是常数，但若选取了几份不同常数，它们仍须在交集上一致才能粘成一个函数。高次形式也有同样的相容问题，只是交集上的差需要用更低次的原形式继续修正。

第一卷定理14.2.2给出了星形开集上的具体算子。若 \(U\subset\mathbb R^n\) 关于 \(0\) 星形、\(k>0\)，则
\[
	(K\omega)_x(v_1,\ldots,v_{k-1})
	=\int_0^1t^{k-1}\omega_{tx}(x,v_1,\ldots,v_{k-1})\,\mathrm{d}t
\]
满足 \(\mathrm{d}K+K\mathrm{d}=\operatorname{id}\)；零次相应公式为 \(K\mathrm{d}f=f-f(0)\)。这里的 \(t^{k-1}\) 来自收缩映射对其余 \(k-1\) 个切向量各乘一次 \(t\)。因此每个正次数闭形式局部有原形式，闭函数局部为常数。以 \(\underline{\mathbb R}\) 表示局部常值实函数层，得到层的正合序列
\begin{equation}\label{g2:eq:de-rham-resolution}
	0\longrightarrow\underline{\mathbb R}\longrightarrow\Omega^0
	\xrightarrow{\mathrm{d}}\Omega^1\xrightarrow{\mathrm{d}}\cdots
	\xrightarrow{\mathrm{d}}\Omega^n\longrightarrow0.
\end{equation}
正合性是在茎上检查的：给定一点的闭形式芽，先缩小到坐标星形邻域再应用上式，并不要求原形式在最初的整个开集存在。每个 \(\Omega^q\) 是光滑函数层的模，单位分解的乘法使它成为细层；\cref{g2:thm:acyclic-resolution}遂给出
\begin{equation}\label{g2:eq:de-rham-sheaf-comparison}
	H^k(M;\underline{\mathbb R})\cong H_{\mathrm{dR}}^k(M).
\end{equation}
这一式比较了层与形式，接下来还需将它与奇异上链的积分映射对应起来。

\begin{definition}[光滑好覆盖]\label{g2:def:good-cover}
	光滑流形上的开覆盖 \(\mathfrak U=(U_i)\) 若满足：每个非空有限交 \(U_{i_0\cdots i_p}\) 都能在自身内光滑收缩到一点，则称为\term[guanghuahaofugai]{光滑好覆盖}[smooth good cover]。本节用形式作收缩时，所说的好覆盖均指这一版本，它比只要求连续可缩的定义多了光滑性条件。
\end{definition}

对这种交集，第一卷定理14.2.1的同伦公式给出正次数 de Rham 上同调为零，零次为 \(\mathbb R\)；光滑棱柱同伦给光滑奇异链相同的结论。因而局部比较只剩常数函数在点上取值，积分在零次正是恒等。

\ProseParagraphBreak

\(\mathbb R^n\) 的任意开集都有可数的这种覆盖：取闭包仍在该开集内、端点为有理数的全部开矩形。有限交若非空仍为凸开矩形，线段收缩留在其中。圆周也可取三个略长于三分之一周的开弧，使每两个相交于一段开弧、三重交为空。这份覆盖的 Čech 复形是三角形的三条边而没有内部面，故 \(\check H^1\cong\mathbb R\)。若只用两个长开弧覆盖圆周，其交有两段，常值层在交集上有两份独立的常数，这份覆盖不满足本定义。

\ProseParagraphBreak

下面先在有好覆盖的区域证明比较，再处理一般流形。一般情形用坐标开集及其交，不要求这些交都可缩；交集可以有洞，只需已经知道其积分比较为同构。这种安排把局部计算与覆盖的几何形状分开。

\subsection{双复形的两种收缩}
\label{g2:subsec:0703:04}
固定可数开覆盖 \(\mathfrak U=(U_i)\)。用上一节的交替指标约定，考虑两个第一象限双复形
\[
	A^{p,q}=\prod_{i_0<\cdots<i_p}\Omega^q(U_{i_0\cdots i_p}),
	\qquad B^{p,q}=\prod_{i_0<\cdots<i_p}C_\infty^q(U_{i_0\cdots i_p};\mathbb R).
\]
横向微分是覆盖的交替限制 \(\delta_{\mathrm C}\)，纵向分别是 \(\mathrm d\) 与奇异上边界 \(\delta_{\mathrm s}\)。两种原始微分交换；总微分必须写成
\[
	\mathrm D_A=\delta_{\mathrm C}+(-1)^p\mathrm d,
	\qquad \mathrm D_B=\delta_{\mathrm C}+(-1)^p\delta_{\mathrm s}
	\quad\text{在双次数 }(p,q)\text{ 上}.
\]
这样两项交叉合成相消，\(\mathrm D^2=0\)。总次数 \(k\) 只含 \((0,k),\ldots,(k,0)\) 这些有限项；无限覆盖的乘积位于每一项内部，不会产生沿总对角线求无限和的问题。

\begin{lemma}[两种增广行的正合性]\label{g2:lem:de-rham-augmented-rows}
	形式的横行增广为 \(\Omega^q(M)\to A^{0,q}\to A^{1,q}\to\cdots\)。奇异上链的横行增广为
	\[
		C_{\infty,\mathfrak U}^q(M;\mathbb R)\longrightarrow B^{0,q}\longrightarrow B^{1,q}\longrightarrow\cdots,
	\]
	其中 \(C_{\infty,\mathfrak U}^q\) 是由像完全落在某个 \(U_i\) 内的光滑单形生成的小链复形的线性对偶。两条增广行均正合，且两种增广均在总上同调上给出同构。
\end{lemma}
\begin{proof}
	对形式选从属于 \(\mathfrak U\) 的光滑单位分解 \((\rho_j)\)。第一卷定理12.3.5保证乘法支集留在相应开集内，且求和局部有限。对交替指标上链，令
	\[
		(ha)_{i_0\cdots i_{p-1}}=\sum_j\rho_j a_{j i_0\cdots i_{p-1}}.
	\]
	每项先在交集定义，再利用 \(\operatorname{supp}\rho_j\subset U_j\) 向其余部分延为零。在边界附近它原本就为零，故延拓仍光滑。展开交替和时，删去 \(j\) 的项为 \((\sum_j\rho_j)a=a\)，其余项两两抵消，得到 \(\delta_{\mathrm C}h+h\delta_{\mathrm C}=\operatorname{id}\)，包括增广位置。

	对奇异上链，不需要给上链乘单位分解。为每条覆盖小单形 \(\sigma\) 选一个指标 \(j(\sigma)\)，使其整个像落在 \(U_{j(\sigma)}\)。定义
	\[
		(ha)_{i_0\cdots i_{p-1}}(\sigma)
		=a_{j(\sigma)i_0\cdots i_{p-1}}(\sigma).
	\]
	只在 \(\sigma\) 完全落在相应交集时求值，故右端有定义。对固定 \(\sigma\)，仍是插入一个首指标的同一交替和消去，得到横向收缩。零次增广所需的上链也按这个逐单形公式定义，不需任何连续性。

	这两个横向算子一般不与纵向微分交换；我们没有把它们直接称为总复形的收缩。\cref{g2:lem:first-quadrant-comparison}按有限总对角线逐项消去多余分量，正好把横行正合性转化为两个增广的总上同调同构。
\end{proof}

小链与全部光滑链也计算相同的上同调。\cref{g2:prop:subdivision-chain-map,g2:prop:subdivision-chain-homotopy,g2:thm:small-chain-isomorphism}的构造只把单形与仿射细分、仿射锥和棱柱复合，因而保持光滑性。每条有限链细分充分多次后落入覆盖小链；圈及填充链都可如此处理。故小链包含诱导同调同构，实数域上的线性对偶再给出上同调同构。

\begin{theorem}[de Rham 定理]\label{g2:thm:de-rham}
	设 \(M\) 为有限维、Hausdorff、第二可数、无边界的光滑流形。积分比较
	\[
		\mathcal I_M:H_{\mathrm{dR}}^k(M)\xrightarrow{\ \cong\ }H^k(M;\mathbb R)
	\]
	在每个次数为同构，并与光滑映射的拉回相容。它与\eqref{g2:eq:de-rham-sheaf-comparison}一道识别形式、常值实数层与奇异上同调；覆盖极限的 Čech 上同调也自然等于这些群。对好覆盖，单份覆盖的常值 Čech 复形已经计算它们，比较由常数包含和积分给出。
\end{theorem}
\begin{proof}
	证明的目标不只是找一个抽象同构，而是核对刚才的积分映射。逐个交集上的积分组成双复形映射 \(I:A^{\bullet,\bullet}\to B^{\bullet,\bullet}\)，因为它与限制及两种纵向微分相容。两条增广在积分下也交换，得到
	\[
		\begin{tikzcd}[column sep=large]
			\Omega^\bullet(M) \arrow[r] \arrow[d,"I"] & \operatorname{Tot} A \arrow[d,"I"] \\
			C_{\infty,\mathfrak U}^\bullet(M;\mathbb R) \arrow[r] & \operatorname{Tot} B .
		\end{tikzcd}
	\]
	上下两个横箭头已在上同调上为同构。

	若 \(\mathfrak U\) 是好覆盖，各交集上的纵向上同调只在零次为 \(\mathbb R\)。积分把常数函数送到同值的零上圈，因此每一列都诱导同构。第一象限比较引理给出中间两个总复形的同构，再由上图得到积分同构。两种双复形还都接受常值 Čech 复形的包含，这两个包含诱导同构且与积分交换，证明所述 Čech 识别。层的识别与此相容，因为常值层到形式消解的增广正是这里使用的常数包含。

	特别地，\(\mathbb R^n\) 的每个开子集都由前述有理开矩形覆盖，故已经满足结论。现在取一般 \(M\) 的可数坐标开覆盖。任一非空有限交包含于其中一个坐标域，在该坐标下是 \(\mathbb R^n\) 的开子集；上一段已给出它的积分同构。于是逐列积分仍诱导同构，虽然这些列的正次上同调未必消失。再用同一双复形比较及增广图，得到一般 \(M\) 的结论。光滑映射的自然性由\cref{g2:prop:integration-naturality}给出。流形满足 Hausdorff、第二可数、局部紧条件，故\cref{g2:thm:cech-sheaf-comparison}还将覆盖极限的 Čech 上同调识别为常值层上同调；这些比较都由同一个常数增广诱导，因而相容。
\end{proof}

在好覆盖上，可以把这段证明理解为两种计算次序：先用单位分解消去覆盖方向，留下整体形式；或者先在可缩交集上求原形式，留下常数及其 Čech 相容条件。一般坐标覆盖则需要保留交集本身的上同调。积分构造与从 Euclid 开集推广到一般流形的思路可对照\cite[第18章，Smooth Singular Homology 与 The de Rham Theorem]{Lee2012Smooth}。

\subsection{乘法和支集版本的比较}
\label{g2:subsec:0703:05}
楔积与杯积的比较要经过乘积空间：先在乘积链上积分，再沿对角映射拉回，就能证明积分同构保持上同调乘法。在单条单形上，两种公式通常不同：楔积在整个单形上积分，Alexander--Whitney 杯积则在前面、后面分别取值。下面的证明将这种链级差别与上同调上的相容区分开来。

\begin{proposition}[乘法的相容性]\label{g2:prop:de-rham-products}
	在\cref{g2:thm:de-rham}的条件下，\(\mathcal I_M\) 保持单位和乘法：
	\[
		\mathcal I_M([\omega]\wedge[\eta])
		=\mathcal I_M[\omega]\smile\mathcal I_M[\eta].
	\]
	对两个这样的流形，积分也把形式的外积 \(\pi_M^*\omega\wedge\pi_N^*\eta\) 送到上同调外积。
\end{proposition}
\begin{proof}
	设 \(\deg\omega=p\)、\(\deg\eta=q\)，两形式闭。\cref{g2:thm:kunneth}在实数域上没有 \(\operatorname{Tor}\) 项，说明 \(H_{p+q}(M\times N;\mathbb R)\) 由各次数的乘积圈类生成。\cref{g2:thm:uct-cohomology}的域版本又说明，一个上同调类由它在全部同调类上的值确定。故只需在 \([z]\times[w]\) 上比较两类，其中 \(z,w\) 可由光滑比较取为光滑圈。

	若 \(\deg z=p\)、\(\deg w=q\)，第4章的乘积链是把有向 \(\Delta^p\times\Delta^q\) 剖成 shuffle 单形后逐项求和。内部公共面相消，定向取第一因子在前。逐片积分再用有限可加性及连续函数在紧乘积区域上的迭代积分公式，得到
	\[
		\int_{z\times w}\pi_M^*\omega\wedge\pi_N^*\eta
		=\left(\int_z\omega\right)\left(\int_w\eta\right).
	\]
	若 \(\deg z=a\)、\(\deg w=b\)、\(a+b=p+q\) 而 \((a,b)\ne(p,q)\)，则 \(a<p\) 或 \(b<q\)，对应因子上的拉回形式为零。上同调外积的值也分别由\cref{g2:prop:cross-cup-evaluation}给出相同的乘积或零。因此外积比较成立。

	最后取 \(N=M\)，沿光滑对角映射 \(\Delta(x)=(x,x)\) 拉回。形式外积拉回为楔积，奇异上同调外积拉回为杯积，自然性给出结论。零次常函数 \(1\) 的积分是常值零上圈 \(1\)，所以单位也保持。这里使用的是同调的域上 Künneth 同构，没有要求上同调外积的张量映射在无限维情形仍为同构。
\end{proof}

紧支比较还要控制原形式的支集：一个闭紧支形式即使有普通原形式，该原形式也可能在无穷远处保持非零。因此，要在紧支复形中判断恰当性，需同时构造紧支的原形式。定义
\[
	\Omega_c^k(M)=\{\omega\in\Omega^k(M):\operatorname{supp}\omega\text{ 紧}\},
	\qquad H_{c,\mathrm{dR}}^k(M)=H^k(\Omega_c^\bullet(M),\mathrm d).
\]
若 \(\omega\) 支于紧集 \(K\)，则 \(I\omega\) 在完全位于 \(M\setminus K\) 的单形上为零，因而是相对上链。第6章的定义给出自然映射
\[
	\mathcal I_{c,M}:H_{c,\mathrm{dR}}^k(M)\longrightarrow
	H_c^k(M;\mathbb R)=\varinjlim_K H^k(M,M\setminus K;\mathbb R).
\]

\begin{theorem}[相对及紧支比较]\label{g2:thm:de-rham-support}
	设 \(M\) 满足\cref{g2:thm:de-rham}的条件，\(V\subset M\) 开。复形
	\[
		A^k(M,V)=\Omega^k(M)\oplus\Omega^{k-1}(V),\qquad
		\mathrm D(\omega,\eta)=(\mathrm d\omega,\omega|_V-\mathrm d\eta)
	\]
	约定负次形式为零，其上同调经积分自然同构于 \(H^k(M,V;\mathbb R)\)。此外，\(\mathcal I_{c,M}\) 在每个次数为同构；它与忘记支集、开嵌入的扩零及固有光滑映射的拉回相容，并把紧支楔积配对送到相应杯积配对。
\end{theorem}
\begin{proof}
	先解释相对复形。一个圈由闭形式 \(\omega\) 及 \(V\) 上满足 \(\omega|_V=\mathrm d\eta\) 的原形式组成；相对条件记录的是这一份原形式数据。直接计算 \(\mathrm D^2=0\)。由投影到 \(\Omega^\bullet(M)\) 得到的短正合列，产生与空间对完全同形的长正合列。

	把形式换成光滑奇异上链，得到相同的复形。上链限制 \(C_\infty^*(M)\to C_\infty^*(V)\) 逐次满射，因为可给其余单形补零；因此相对上链的核包含 \(\varphi\mapsto(\varphi,0)\) 诱导此复形的上同调同构。也可逐项验证：对一个圈 \((a,b)\)，延拓 \(b\) 为 \(\widetilde b\)，减去余边界 \(\mathrm D(\widetilde b,0)\) 后，第二项为零，第一项落在限制核内；相同操作处理边界。积分与两条长正合列交换，对 \(M,V\) 用普通 de Rham 定理及五引理，得到相对比较。光滑与连续相对模型的比较则由绝对比较和对的长正合列得到。

	现在令 \(V=M\setminus K\)，并使紧集 \(K\) 增大。证明自然包含
	\[
		\Omega_c^\bullet(M)\longrightarrow\varinjlim_K A^\bullet(M,M\setminus K),
		\qquad\omega\longmapsto(\omega,0)
	\]
	在上同调上为同构。给定右侧圈 \((\omega,\eta)\)，取光滑函数 \(\rho\)，使其在 \(K\) 的一个邻域为零，在某个更大紧集 \(L\) 外为 \(1\)。第一卷推论12.3.6的闭集截断可作此选择。把 \(\rho\eta\) 延为 \(M\) 上的形式 \(\widetilde\eta\)；它在 \(K\) 附近为零，所以延拓光滑。在更大紧集 \(L\) 对应的相对复形中，减去 \(\mathrm D(\widetilde\eta,0)\) 后得到
	\[
		(\omega-\mathrm d\widetilde\eta,0),
	\]
	且第一项在 \(M\setminus L\) 为零。这证明上同调满射。零次没有 \(\eta\)，同一结论直接由 \(\omega|_{M\setminus K}=0\) 得到。

	为证单射，设紧支闭形式 \(\omega\) 在极限中为边界。在足够大紧集 \(K\) 的阶段，有 \(\omega=\mathrm d\beta\)、\(\beta|_{M\setminus K}=\mathrm d\gamma\)。用相同截断延拓 \(\rho\gamma\)，则 \(\beta-\mathrm d(\rho\gamma)\) 是紧支形式，其外微分仍为 \(\omega\)。当 \(\deg\omega=1\) 时 \(\gamma\) 的次数为负，条件已经说 \(\beta\) 在补集为零；零次没有非零边界。故各次均得单射。向量空间的滤过直接极限保持正合：圈、边界及两个代表相等都由有限个等式见证，故在某个更晚阶段已经成立。于是可交换这里的上同调与直接极限，结合相对比较得到 \(\mathcal I_{c,M}\) 同构。

	忘记支集与固有拉回在积分链级相容；固有性保证紧集的逆像仍紧。开嵌入扩零则在包含支集的相对对上由切除实现，与形式的扩零相同。乘法证明也保留支集：若第一形式支于 \(K\)，就在空间对 \((M,M\setminus K)\) 上作外积比较。乘积链映射及其同伦保持补集子链，故第4章的域上 Künneth 论证下降到相对商；上述迭代积分计算仍成立。沿对角拉回后，支集回到 \(K\)。若两形式都紧支，所得支集包含在两支集的交中，取直接极限便得到紧支杯积配对。
\end{proof}

取 \(f\in C_c^\infty(\mathbb R)\)、\(\int_{\mathbb R}f(x)\,\mathrm dx=1\)。形式 \(f(x)\,\mathrm dx\) 有普通原函数 \(F(x)=\int_{-\infty}^x f(t)\,\mathrm dt\)，但 \(F\) 在正向远端为 \(1\)，不可能紧支；任何其他原函数只差常数，也不能在两端都为零。因此该形式在 \(H_{c,\mathrm{dR}}^1(\mathbb R)\) 非零，在普通 \(H_{\mathrm{dR}}^1(\mathbb R)\) 却为零。这正对应第6章 \(H_c^1(\mathbb R;\mathbb R)=\mathbb R\) 与 \(H^1(\mathbb R;\mathbb R)=0\) 的区别。

% !TeX root = ../../Book2.tex
% 纲要标题：### 7.4 de Rham 上同调的计算与局部系数
% 纲要位置：几何学计划纲要.md，第 1290 行。
% 正文状态：初稿；尚未正式审核。

% 纲要内容开始
% 
% * 周期与整上同调
% * 环面、球面及射影空间计算
% * 局部系统值形式的平坦联络模型
% 
% ---
% 
% 纲要内容结束

\section{de Rham 上同调的计算与局部系数}
\label{g2:sec:0704}

积分比较把一个闭形式变成对圈的测量，但它只给出实系数信息。一个周期是整数，究竟是对形式本身的限制，还是对它所代表的类的限制？如果某个整数类有有限阶，它还能由一个非零闭形式表示吗？先澄清这些问题，再把第4章的环计算落实为具体微分形式。最后，我们把实数换成沿道路运输的向量，说明第5章的局部系统怎样进入微分方程。

本节的流形均为有限维、Hausdorff、第二可数的无边界光滑流形；紧性只在明确提出时使用。积分所用链为上一节的光滑奇异链，系数为整数或实数。由其与连续奇异链的比较，它们计算同样的同调群。

\subsection{周期与整上同调}
\label{g2:subsec:0704:01}
设 \(\omega\in\Omega^k(M)\) 满足 \(\mathrm d\omega=0\)。对光滑整数 \(k\) 圈 \(z=\sum_i n_i\sigma_i\)，定义其\term[zhouqi]{周期}[period]为
\[
	\operatorname{Per}_{\omega}([z])=\int_z\omega
	\coloneqq\sum_i n_i\int_{\Delta^k}\sigma_i^*\omega.
\]
若 \(z-z'=\partial c\)，Stokes 公式给出 \(\int_{z-z'}\omega=\int_c\mathrm d\omega=0\)；若以 \(\omega+\mathrm d\eta\) 替换 \(\omega\)，在圈上的积分也不变。因此周期定义了群同态 \(H_k(M;\mathbb Z)\to\mathbb R\)，且只依赖 \([\omega]\in H^k_{\mathrm{dR}}(M)\)。当 \(k=0\) 时，闭零形式是局部常函数，它在点所代表的零圈上的周期就是函数值。

整数链上的积分仍取实数值：例如在 \(\mathbb R/\mathbb Z\) 上，\(a\,\mathrm dt\) 的一周积分是任意指定的 \(a\in\mathbb R\)。另一方面，积分在挠类上恒为零：若 \(r[z]=0\)，则 \(r\int_z\omega=0\)，故 \(\int_z\omega=0\)。整周期条件所限制的是这些实数取值是否全部属于 \(\mathbb Z\)。

\begin{proposition}[整周期判据]\label{g2:prop:integral-periods}
	设 \(M\) 为无边界光滑流形，\(k\geq0\)，并且 \(H_k(M;\mathbb Z)\)、\(H_{k-1}(M;\mathbb Z)\) 有限生成，其中 \(H_{-1}=0\)。对闭形式 \(\omega\in\Omega^k(M)\)，下列条件等价：
	\begin{enumerate}
		\item 对每个光滑整数 \(k\) 圈 \(z\)，都有 \(\int_z\omega\in\mathbb Z\)。
		\item 在\cref{g2:thm:de-rham}的积分同构下，\([\omega]\) 属于系数映射 \(H^k(M;\mathbb Z)\to H^k(M;\mathbb R)\) 的像。
	\end{enumerate}
	该像同构于 \(H^k(M;\mathbb Z)/\operatorname{Tor}H^k(M;\mathbb Z)\)，在 \(H^k(M;\mathbb R)\) 中是一个满格。特别地，此判据适用于紧光滑流形。
\end{proposition}
\begin{proof}
	先恢复万有系数定理在这里的作用。\cref{g2:thm:uct-cohomology}的评价映射给出交换图
	\[
		\begin{tikzcd}[column sep=large]
			H^k(M;\mathbb Z) \arrow[r,two heads,"\rho_{\mathbb Z}"] \arrow[d] & \operatorname{Hom}(H_k(M;\mathbb Z),\mathbb Z) \arrow[d] \\
			H^k(M;\mathbb R) \arrow[r,"\rho_{\mathbb R}"'] & \operatorname{Hom}(H_k(M;\mathbb Z),\mathbb R).
		\end{tikzcd}
	\]
	下方箭头是同构：有限生成群分解为自由群与有限循环群的直和，而 \(\operatorname{Ext}^1(\mathbb Z,\mathbb R)=0\)、\(\operatorname{Ext}^1(\mathbb Z/r,\mathbb R)=\mathbb R/r\mathbb R=0\)，后一个公式来自\cref{g2:prop:cyclic-ext-tor}。因此实上同调类完全由其在整数圈上的取值确定。

	积分比较将 \(\rho_{\mathbb R}[\omega]\) 识别为刚定义的周期同态。如果它的值全部为整数，就得到一个同态 \(H_k(M;\mathbb Z)\to\mathbb Z\)；上方箭头满射，故这个同态提升为整数上同调类，交换性和下方箭头的单射性说明其实像正是 \([\omega]\)。反过来，整数上同调类在整数圈上当然取整数值，交换图给出所需结论。

	写 \(H_k(M;\mathbb Z)\cong\mathbb Z^b\oplus T\)，其中 \(T\) 有限。右上角是 \(\mathbb Z^b\)，右下角是 \(\mathbb R^b\)，右侧映射是通常的整数格包含。左上角评价的核为 \(\operatorname{Ext}^1(H_{k-1}(M;\mathbb Z),\mathbb Z)\)，由同样的有限生成分解可知它是有限群；评价的目标无挠，所以这个核恰是全部挠子群。由此得到最后的断言。紧流形的整数同调有限生成已由\cref{g2:lem:compact-manifold-finite-homology}证明。
\end{proof}

具有上述性质的闭形式称为具有\term[zhengzhouqi]{整周期}[integral periods]。这个条件既不要求形式在坐标中的系数为整数，也不指定唯一的整数上同调提升。后者可以相差一个挠类，而所有周期仍完全相同。实上同调是向量空间，整周期类只是其中的格；任意实倍数通常会离开这个格。

还有一个较弱的条件应单独辨认：\(\omega\) 为闭形式只说明积分对同调不变，并不说明任何周期为零。恰当形式的周期全部为零；在上述有限生成条件下，反方向由积分同构及评价同构成立，即所有周期为零恰好等价于 \(\omega\) 恰当。对 \(S^1=\mathbb R/\mathbb Z\)，\(\mathrm dt\) 闭且有整周期，但其一周积分为 \(1\)，所以不是恰当形式。

\subsection{环面、球面及射影空间计算}
\label{g2:subsec:0704:02}
在 \(T^m=\mathbb R^m/\mathbb Z^m\) 上，坐标函数 \(t_i\) 不下降为实值函数，其微分却下降为整体一形式 \(\theta_i\)：整数平移改变函数一个常数，不改变微分。令 \(\gamma_j\) 为正向的第 \(j\) 条坐标圆周，则
\[
	\int_{\gamma_j}\theta_i=\delta_{ij}.
\]
所以积分同构将 \([\theta_i]\) 送到\cref{g2:prop:torus-cohomology-ring}中的第 \(i\) 个一次生成元的实像。该同构保持乘法，因而
\begin{equation}\label{g2:eq:torus-de-rham-ring}
	H^*_{\mathrm{dR}}(T^m)
	=\Lambda_{\mathbb R}([\theta_1],\ldots,[\theta_m]).
\end{equation}
具体说，每个闭 \(k\) 形式都与唯一一个常系数组合 \(\sum_{i_1<\cdots<i_k}a_{i_1\cdots i_k}\theta_{i_1}\wedge\cdots\wedge\theta_{i_k}\) 相差恰当形式；这些系数就是它在相应坐标子环面上的积分。因而该类具有整周期，当且仅当所有 \(a_{i_1\cdots i_k}\) 为整数。这些常系数代表由积分比较和环面的杯积计算共同确定。

对 \(n\geq1\)，在单位球面 \(S^n\subset\mathbb R^{n+1}\) 上取向外定向，并令
\[
	\nu=\frac{1}{\operatorname{vol}(S^n)}
	\left.\sum_{j=0}^n(-1)^j x_j\,\mathrm dx_0\wedge\cdots\wedge\widehat{\mathrm dx_j}\wedge\cdots\wedge\mathrm dx_n\right|_{S^n}.
\]
分母是通常球面面积，故 \(\int_{S^n}\nu=1\)。顶次形式自动闭；它不是恰当形式，否则 Stokes 公式会使这个积分为零。球面的同调计算与积分比较给出
\[
	H^k_{\mathrm{dR}}(S^n)=
	\begin{cases}
		\mathbb R[1],&k=0,\\
		\mathbb R[\nu],&k=n,\\
		0,&\text{其余次数}.
	\end{cases}
\]
全部正次乘积为零。任意顶次形式 \(\omega\) 都有 \([\omega]=(\int_{S^n}\omega)[\nu]\)；这给出一个无需显式求原函数的恰当性判据。

复射影空间需要一个二形式，而不只是顶次体积形式。下面从单位球面构造它，使归一化可以直接检查。记 \(q:S^{2m+1}\to\mathbb CP^m\) 为 \(z\mapsto[z]\)，其中 \(m\geq1\)、\(z_j=x_j+\mathrm i y_j\)，并在球面上取
\[
	\alpha=\frac{1}{2\pi}\sum_{j=0}^m(x_j\,\mathrm dy_j-y_j\,\mathrm dx_j).
\]
圆周作用为 \(z\mapsto e^{\mathrm i\vartheta}z\)，正向参数 \(0\leq\vartheta\leq2\pi\) 的纤维上有 \(\alpha=\mathrm d\vartheta/(2\pi)\)。在仿射图中，把某个非零齐次坐标选为正实数再归一化，得到光滑局部截面 \(s\)。两条这样的截面在足够小的交叠上满足 \(s'=e^{\mathrm i f}s\)，其中 \(f\) 是实值光滑函数；直接代入给出
\[
	(s')^*\alpha=s^*\alpha+\frac{\mathrm df}{2\pi}.
\]
因此 \(\mathrm d(s^*\alpha)\) 在交叠上一致，粘成 \(\mathbb CP^m\) 上的闭二形式 \(\omega\)。这也说明 \(q^*\omega=\mathrm d\alpha\)：在局部乘积坐标中，\(\alpha\) 等于截面上的形式加上纤维角度的微分，而后者的外微分为零。\(\mathrm d\alpha\) 在球面上恰当，并不意味着下降后的 \(\omega\) 也恰当；原始的一形式 \(\alpha\) 在纤维方向不为零，不能下降到射影空间。

在坐标直线 \(\mathbb CP^1\) 的图 \([1:z]\)、\(z=x+\mathrm i y\) 中，可取 \(s(z)=(1,z)/\sqrt{1+|z|^2}\)。于是
\[
	s^*\alpha=\frac{x\,\mathrm dy-y\,\mathrm dx}{2\pi(1+x^2+y^2)},
	\qquad \omega=\frac{\mathrm dx\wedge\mathrm dy}{\pi(1+x^2+y^2)^2}.
\]
取复坐标给出的定向，利用极坐标积分得到
\[
	\int_{\mathbb CP^1}\omega
	=\frac1\pi\int_0^{2\pi}\int_0^\infty\frac{r}{(1+r^2)^2}\,\mathrm dr\,\mathrm d\vartheta=1.
\]
被省略的点没有面积，且形式在另一张图上光滑，故这里的广义积分等于紧流形上的积分。\cref{g2:thm:complex-projective-ring}说明 \(H^2(\mathbb CP^m;\mathbb Z)\to H^2(\mathbb CP^1;\mathbb Z)\) 是同构，所以 \([\omega]\) 是在正向坐标直线上取值 \(1\) 的整数生成元的实像。乘法相容性遂给出
\begin{equation}\label{g2:eq:projective-de-rham-ring}
	H^*_{\mathrm{dR}}(\mathbb CP^m)
	\cong\mathbb R[u]/(u^{m+1}),\qquad u=[\omega],\quad |u|=2.
\end{equation}
特别地，\(1,\omega,\ldots,\omega^m\) 的类分别给出全部偶数次的一维上同调，奇数次为零。非零的楔幂在这里由第4章已经证明的整数杯积计算保证。

实射影空间则清楚显示实系数丢失了什么。\cref{g2:prop:projective-integer-homology}与万有系数定理给出，当 \(m\geq1\) 时
\[
	H^k(\mathbb RP^m;\mathbb Z)=
	\begin{cases}
		\mathbb Z,&k=0,\text{ 或 }k=m\text{ 为奇数},\\
		\mathbb Z/2,&k\text{ 为正偶数且 }k\leq m,\\
		0,&\text{其余次数}.
	\end{cases}
\]
这里偶数次的 \(\mathbb Z/2\) 来自前一次同调的 Ext 项。换成实系数以后，这些项全部消失，故
\[
	H^k_{\mathrm{dR}}(\mathbb RP^m)=
	\begin{cases}
		\mathbb R,&k=0,\text{ 或 }k=m\text{ 为奇数},\\
		0,&\text{其余次数}.
	\end{cases}
\]
例如 \(H^2(\mathbb RP^2;\mathbb Z)=\mathbb Z/2\)，却没有非零的二次 de Rham 类。若 \(m\) 为奇数，球面对径映射保持定向，归一化球面体积形式 \(\nu\) 下降为 \(\bar\nu\)，且二重覆盖给出 \(\int_{\mathbb RP^m}\bar\nu=1/2\)；因此具有整周期的正向顶次生成元由 \(2\bar\nu\) 表示。当 \(m\) 为偶数时，对径映射反转定向，这个体积形式不能作为普通实值形式下降。

\subsection{局部系统值形式的平坦联络模型}
\label{g2:subsec:0704:03}
第5章的局部系统沿道路比较不同位置的系数。若希望用微分形式作同样的计算，首先要知道怎样微分一个取值于不同纤维的对象。设 \(E\to M\) 是有限秩实光滑向量丛，取联络 \(\nabla\)，并记
\[
	\Omega^k(M;E)=\Gamma(\Lambda^k\mathrm T^*M\otimes E).
\]
沿用第一卷第15章的行标架 \(e=(e_1,\ldots,e_r)\) 与列分量约定：若 \(\nabla e=eA\)，则 \(\nabla(eu)=e(\mathrm du+Au)\)。形式上的\term[waixiebianweifen]{外协变微分}[exterior covariant derivative]定义为
\begin{equation}\label{g2:eq:flat-valued-differential}
	\mathrm d_\nabla(e\eta)=e(\mathrm d\eta+A\wedge\eta).
\end{equation}
换标架 \(e'=eg\) 时，\(A'=g^{-1}Ag+g^{-1}\mathrm dg\)、\(\eta'=g^{-1}\eta\)，代入可知右侧表示同一个整体形式。因此这一定义不依赖标架。展开两次求导并使用 \(\mathrm d(A\wedge\eta)=\mathrm dA\wedge\eta-A\wedge\mathrm d\eta\)，得到
\begin{equation}\label{g2:eq:flat-curvature-square}
	\mathrm d_\nabla^2(e\eta)=e(F\wedge\eta),\qquad
	F=\mathrm dA+A\wedge A.
\end{equation}
这正是第一卷定理15.3.5的曲率公式。只有曲率为零的\term[pingtanlianluo]{平坦联络}[flat connection]才能给出所需上链复形；要求平方在零形式上都为零，就已经强迫 \(F=0\)。

平行标架使局部联络矩阵为零，于是向量丛值形式的外协变微分就化为逐分量的普通外微分。第一卷定理15.5.2已经证明，\(F=0\) 等价于每点附近存在这样的标架。在坐标盒中写 \(A=\sum_i A_i\,\mathrm dx^i\)，其构造先解带参数的线性常微分方程
\[
	\partial_1g=-A_1g,\qquad g(0,x^2,\ldots,x^n)=I.
\]
解对参数光滑；同时解 \(\partial_1h=hA_1\)、\(h(0,x')=I\)，可得 \(hg=I\)，故换标架合法。换后 \(B_1=0\)，而零曲率给出 \(0=F_{1j}'=\partial_1B_j\)，所以其余系数只依赖 \(x^2,\ldots,x^n\)。在这些变量上归纳消去剩余系数，就得到 \(A=0\) 的标架。这里用到的参数常微分方程定理及光滑依赖已在第一卷第13.1节证明；每一步可缩小坐标盒，以取得共同定义域。反过来，平行标架中 \(A=0\) 显然使曲率为零。

两个平行标架满足 \(e'=eg\) 时，\(0=g^{-1}\mathrm dg\)，所以转移矩阵局部常值。在这样的标架中保持向量分量不变，便能沿道路运输。把一个道路同伦的参数方形细分到每个小格都落在一个平行标架域，格子内两条边路的运输相同，公共边逐一抵消，证明运输只依赖定端点同伦类；这就是第一卷定理15.5.4的论证。因此平行截面组成有限维实\term[jubu xitong]{局部系统}[local system] \(\mathcal L\)。其纤维可识别为 \(E_x\)，运输即平行移动。反过来，把有限秩实局部系统的局部常转移矩阵用于粘合 \(U_i\times\mathbb R^r\)，便得到光滑向量丛；因为转移矩阵的微分为零，各片的普通微分粘成平坦联络。这两个构造在相同局部平行标架中互相恢复，也保持纤维线性且与运输相容的映射。

\begin{theorem}[局部系统的 de Rham 比较]\label{g2:thm:twisted-de-rham}
	设 \(M\) 为无边界光滑流形，\(E\to M\) 为有限秩实光滑向量丛，\(\nabla\) 为平坦联络，\(\mathcal L\) 为其平行截面局部系统。则运输后的积分诱导自然同构
	\[
		H^k\bigl(\Omega^*(M;E),\mathrm d_\nabla\bigr)
		\cong H^k(M;\mathcal L).
	\]
	对常值系统及其平凡平坦丛，此同构就是\cref{g2:thm:de-rham}的积分同构。
\end{theorem}
\begin{proof}
	先说明积分取值在哪里。对光滑单形 \(\sigma:\Delta^k\to M\)，将 \(\sigma^*E\) 的各纤维沿单形内道路运输到首顶点 \(v_0\)。单形可缩，平坦运输使这个比较与道路无关。于是 \(\sigma^*\eta\) 成为取值于固定向量空间 \(E_{\sigma(v_0)}\) 的形式，可逐分量积分；记所得元素为 \(I_\nabla\eta(\sigma)\)。这种描述没有选择基，因而积分与平行标架的改变相容。

	在首顶点平行标架中，\(\mathrm d_\nabla\) 变成普通外微分，Stokes 公式逐分量成立。边界中去掉 \(v_0\) 的面以 \(v_1\) 为首顶点，必须先把它的积分沿边 \([v_0,v_1]\) 反向运回 \(v_0\)；其余各面首顶点仍是 \(v_0\)。因此
	\begin{align*}
		I_\nabla(\mathrm d_\nabla\eta)(\sigma)
		&=T_{01}^{-1}I_\nabla\eta(\sigma|[v_1,\ldots,v_{k+1}])\\
		&\quad+\sum_{i=1}^{k+1}(-1)^iI_\nabla\eta(\sigma|[v_0,\ldots,\widehat v_i,\ldots,v_{k+1}]).
	\end{align*}
	这恰好是\cref{g2:def:twisted-cochains}的余边界公式，故积分是上链映射。上一节的相对光滑逼近比较也适用于这个局部系统：每个单形同伦的运输把新旧首顶点纤维对应起来，棱柱公式中的公共面仍以相反符号和相同运输抵消。因此光滑局部系数上链与连续局部系数上链给出同一上同调。

	在平行标架坐标小球内，形式复形在正次数的局部恰当性逐分量成立，零次核就是 \(\mathcal L\)。因此层复形
	\[
		0\longrightarrow\mathcal L\longrightarrow\Omega^0(E)
		\xrightarrow{\mathrm d_\nabla}\Omega^1(E)
		\xrightarrow{\mathrm d_\nabla}\cdots
	\]
	正合。形式层可以乘光滑单位分解，故由本章的细层无上同调结论及\cref{g2:thm:acyclic-resolution}，其整体截面计算 \(\mathcal L\) 的层上同调。还须把这项计算与奇异上同调通过指定的积分联系起来。

	取由平行标架坐标域 \(U_i\) 组成的开覆盖。每个非空有限交 \(U_{i_0\cdots i_p}\) 是第一张坐标图中的 Euclid 开集，且第一张图的平行标架在交集上仍然平凡化 \(\mathcal L\)。该交集未必可缩，但\cref{g2:thm:de-rham}已经对任意 Euclid 开集证明普通积分同构，逐分量应用便得到这个交集上带 \(\mathcal L\) 系数的积分同构。现在将上一节的覆盖双复形换成 \(E\) 值形式与 \(\mathcal L\) 值小单形上链。形式增广行仍由光滑单位分解收缩；上链增广行仍可对每个小单形选择一个容纳它的覆盖指标作收缩，限制不改变其首顶点纤维，因而\cref{g2:lem:de-rham-augmented-rows}的证明保持有效。刚才在有限交上的积分同构说明每列的比较是同构，\cref{g2:lem:first-quadrant-comparison}遂给出总复形同构，再沿两个增广映射得到所述整体积分同构。

	光滑映射的拉回与单形上的运输、积分都交换；保持联络的丛映射在平行标架中是常矩阵，也与积分交换。这给出所述自然性。
\end{proof}

\begin{example}[圆周上的平坦系数]\label{g2:exm:flat-circle-de-rham}
	设 \(E\to S^1=\mathbb R/\mathbb Z\) 的实际正向环路运输为 \(T\in\operatorname{GL}(V)\)。在区间 \([0,1]\) 上沿路径平行平凡化，截面的列函数满足 \(f(1)=T^{-1}f(0)\)，而不是 \(f(1)=Tf(0)\)：终点的平行标架本身已相对于起点乘了 \(T\)。拉回到实线后，形式写作 \(f(t)\) 或 \(g(t)\,\mathrm dt\)，并满足所有导数的相应准周期条件。

	此时微分为普通导数。零次核是满足 \(T^{-1}v=v\) 的常函数，故为 \(\ker(T-1)\)。对一次形式，令 \(b=\int_0^1g(t)\,\mathrm dt\)。若它是 \(\mathrm df\)，则
	\[
		b=f(1)-f(0)=(T^{-1}-1)f(0).
	\]
	反过来，若 \(b=(T^{-1}-1)v\)，取 \(f(t)=v+\int_0^t g(s)\,\mathrm ds\)。端点值满足准周期关系；一阶以上导数的关系来自 \(g\)，所以 \(f\) 下降为光滑截面。任给 \(b\in V\)，取支集位于 \((0,1)\)、积分为 \(1\) 的光滑函数 \(\chi\)，令 \(g=\chi b\) 并作准周期延拓，就实现这个积分。于是
	\[
		H^0_{\mathrm d_\nabla}(S^1;E)=\ker(T-1),\qquad
		H^1_{\mathrm d_\nabla}(S^1;E)=V/(T^{-1}-1)V=V/(T-1)V.
	\]
	这与\cref{g2:prop:circle-local-coefficients}的胞腔计算相合，并解释了一次余边界为什么含逆运输。若 \(V=\mathbb R\)、\(T=-1\)，两个群均为零；但整数符号系统的一次上同调为 \(\mathbb Z/2\)。把系数几何化为实平坦丛依然不能恢复已经被实数消掉的挠信息。
\end{example}

第5章把纤维的实同调或上同调组织为基底上的局部系统；当基底光滑且这些纤维向量空间有限维时，刚才的构造将它们实现为平坦丛。于是沿基底变化的纤维上同调类可以写成平坦丛截面，\(\mathrm d_\nabla\) 测量它是否局部保持平行。这里要求有限维，且只得到所给局部系统的微分形式模型；一个局部系统仍不包含原几何族的全部粘合信息，这与\secref{g2:sec:0504}的区别一致。
% 纲要条目排为编号小节标题，并进入目录与书签；说明段仍保留于纲要注释中。


\begin{chaptersummary}
	层把开集上的截面及其限制组织在一起；层化允许局部可表示的数据粘合，茎则只保留一点附近的信息。核、余核与正合性可以逐茎检测，整体截面却只自动保持左正合。逐点茎乘积的典范消解把整体提升失败定义为高次层上同调；松弛性允许延拓，细性允许带支集地加权粘合，两者各自给出消失机制。

	\ProseParagraphBreak

	Čech 微分交替删去指标，记录局部数据在各级交叠上的相容关系。细化指标的选择通过链同伦消失。无上同调覆盖可直接计算层上同调；在 Hausdorff、第二可数、局部紧空间中，即使不能一次找到这样的覆盖，也可以通过细化逐次消掉固定总次数中的局部障碍，得到覆盖极限的比较同构。

	\ProseParagraphBreak

	光滑逼近必须保持公共面，否则链边界的抵消会丢失。紧像附近的欧氏模型、角面延拓与相对平滑给出光滑链比较。Stokes 公式随后使积分成为上链映射；覆盖双复形一边通过单位分解恢复整体形式，另一边通过局部积分比较接到奇异上同调。楔积与杯积在上链上一般不相等，但乘积积分及对角拉回证明其在上同调上相容。

	\ProseParagraphBreak

	相对形式复形同时保留一个闭形式及其在子空间上的原形式；紧支比较还须通过截断消去远处的原形式。实形式的整周期可检测整数上同调在实上同调中的像，却不能区分其中的挠类。平坦联络的局部平行标架把局部系统变为向量丛值形式，曲率为零正是外协变微分平方为零的条件；圆周上的准周期函数又把这一模型还原为第5章的单值核与余核计算。
\end{chaptersummary}

\BookExercises

\begin{exercise}\label{g2:ex:skyscraper-sheaf}
	设 \(X\) 为 Hausdorff 空间，\(x\in X\)，\(A\) 为 \(R\) 模。定义 \(\mathcal S_x(A)(U)=A\)（\(x\in U\)）及 \(0\)（\(x\notin U\)），限制在两端均含 \(x\) 时取恒等，其余取零。证明它是层，计算各点的茎及全部层上同调。进一步证明，对任意 \(R\) 模层 \(\mathcal F\)，有自然同构
	\[
		\operatorname{Hom}(\mathcal F,\mathcal S_x(A))\cong\operatorname{Hom}_R(\mathcal F_x,A).
	\]
\end{exercise}
\begin{solution}
	若覆盖的开集之并不含 \(x\)，所有截面都是零。若并含 \(x\)，至少一个覆盖成员含 \(x\)，而所有含 \(x\) 的成员两两交也含 \(x\)，所以相容族必须给出同一个 \(a\in A\)。它唯一粘成并集上的 \(a\)，故这是层。在 \(x\) 的每个邻域上模都为 \(A\)，限制为恒等，故茎为 \(A\)。若 \(y\ne x\)，Hausdorff 性给出不含 \(x\) 的 \(y\) 邻域，其上的模为零，故 \(y\) 处的茎为零。
	所有限制均满射，因此该层松弛。由\cref{g2:thm:flabby-fine-acyclic}，\(H^0(X,\mathcal S_x(A))=A\)，而 \(H^q(X,\mathcal S_x(A))=0\)（\(q>0\)）。
	给定层态射 \(f\)，它在含 \(x\) 的开集上的值与限制相容，遂定义 \(\mathcal F_x\to A\)，\(s_x\mapsto f_U(s)\)。若代表改变，在某个共同小邻域上相等，所得值不变。反之，给定 \(u:\mathcal F_x\to A\)，在含 \(x\) 的开集上设 \(f_U(s)=u(s_x)\)，在其余开集上设为零；它显然与限制交换。这两个构造互逆，并与两边的模或层态射相容，故自然。这样的层称为集中于 \(x\) 的摩天楼层。
\end{solution}

\begin{exercise}\label{g2:ex:exponential-sheaf-circle}
	在 \(S^1\) 上令 \(\mathcal C_{\mathbb R}\) 为连续实函数层，\(\mathcal C_{S^1}\) 为连续圆值函数层，后者采用逐点乘法。证明
	\[
		0\longrightarrow\underline{\mathbb Z}\longrightarrow\mathcal C_{\mathbb R}\xrightarrow{\ f\mapsto e^{2\pi\mathrm i f}\ }\mathcal C_{S^1}\longrightarrow0
	\]
	为阿贝尔群层短正合列。由此计算 \(H^1(S^1,\underline{\mathbb Z})\)，并说明连接同态如何记录连续映射 \(g:S^1\to S^1\) 的绕数。
\end{exercise}
\begin{solution}
	指数映射的核是连续整数值函数，也就是局部常值整数函数。对任意 \(g\) 和任意 \(x\)，可缩小 \(x\) 的邻域，使 \(g\) 的像落在圆的一个真开弧内；这个开弧上有连续角度函数，故指数映射在茎上满射。注意这没有断言每个整体圆值函数都有整体角度。
	连续单位分解把 \(\mathcal C_{\mathbb R}\) 变成细层：端同态为乘以连续权函数。由第一卷第4章的单位分解定理与\cref{g2:thm:flabby-fine-acyclic}，它在圆上无上同调。因此长正合列给出
	\[
		H^1(S^1,\underline{\mathbb Z})\cong
		 C(S^1,S^1)/\{e^{2\pi\mathrm i f}:f\in C(S^1,\mathbb R)\}.
	\]
	把圆参数化为 \(t\mapsto e^{2\pi\mathrm it}\)，\(0\le t\le1\)。第一卷第5章的道路提升定理对覆盖 \(\mathbb R\to S^1\) 给出 \(g(e^{2\pi\mathrm it})=e^{2\pi\mathrm i\theta(t)}\) 的连续提升。端点投影相同，故 \(k=\theta(1)-\theta(0)\in\mathbb Z\)。改变初值只使整个提升加同一个整数，故 \(k\) 不变；相乘的提升可相加，所以绕数是群同态。
	若 \(g=e^{2\pi\mathrm if}\) 有整体角度，则取 \(\theta(t)=f(e^{2\pi\mathrm it})\)，两端相等，绕数为零。反之绕数为零时 \(\theta(0)=\theta(1)\)，所以 \(\theta\) 下降为圆上的连续函数 \(f\)。每个整数 \(k\) 都由 \(g(z)=z^k\) 实现，故上述商群恰为 \(\mathbb Z\)。以该识别作约定，连接类 \(\partial g\) 对应绕数 \(k\)；零类恰好等价于整体角度存在。
\end{solution}

\begin{exercise}\label{g2:ex:twisted-cech-circle}
	设 \(A\) 为 \(R\) 模，\(T:A\to A\) 为自同构。在圆的两弧覆盖 \(U,V\) 上给定局部常值层：两个交集分支的过渡映射分别为 \(1\) 与 \(T\)。在共同坐标中把 Čech 微分写成 \(\delta(a,b)=(b-a,b-Ta)\)。计算该覆盖的零次、一次 Čech 上同调，并分别讨论 \(A=\mathbb Z,T=-1\) 和 \(A=\mathbb R,T=-1\)。
\end{exercise}
\begin{solution}
	零次核满足 \(b=a=Ta\)，所以
	\[
		\check H^0\cong\ker(T-1),\qquad
		\check H^1\cong A/(T-1)A.
	\]
	第二个同构可由 \((u,v)\mapsto[v-u]\) 给出：微分的像映到 \([(1-T)a]=0\)。它满射，因为 \((0,w)\) 映到 \([w]\)。若 \(v-u=(1-T)a\)，取 \(b=u+a\)，则 \(b-a=u\)、\(b-Ta=u+(1-T)a=v\)，故该映射的核恰是微分的像。不存在三重交的项，所以高次 Čech 群为零。
	当 \(A=\mathbb Z,T=-1\) 时，不变元只有零，余不变元为 \(\mathbb Z/2\)。当 \(A=\mathbb R,T=-1\) 时，乘 \(2\) 可逆，零次与一次群都为零。相同的反号运输可以在整数系数下留下挠类，在实系数下完全消失；这不是由覆盖形状单独决定的。本题计算的是指定覆盖复形，认同于层上同调时还须核验有限交上的无上同调条件。
\end{solution}

\begin{exercise}\label{g2:ex:smooth-jet-cokernel}
	设 \(m\ge1\)，在 \(\mathbb R\) 的光滑函数层 \(\mathcal C^{\infty}\) 上考虑乘法态射 \(f\mapsto x^m f\)。证明其核为零，余核同构于集中在 \(0\) 的摩天楼层 \(\mathcal S_0(\mathbb R^m)\)，其中商映射记录 \((f(0),f'(0),\ldots,f^{(m-1)}(0))\)。解释为什么这个余核并不是一个在每点都有 \(m\) 维茎的常值层。
\end{exercise}
\begin{solution}
	若 \(x^m f(x)=0\)，则在所有非零点 \(f=0\)，连续性又给出 \(f(0)=0\)，故核层为零。定义层态射 \(J:\mathcal C^{\infty}\to\mathcal S_0(\mathbb R^m)\)：含 \(0\) 的开集上取所列导数，不含 \(0\) 时取零。任意给定的 \(m\) 个导数值都由一个次数小于 \(m\) 的多项式实现，故 \(J\) 在茎上满射。
	在非零点附近，\(x^m\) 可逆，像层就是全部光滑函数层。在 \(0\) 附近，若 \(f^{(j)}(0)=0\)（\(0\le j<m\)），Taylor 公式的积分余项给出
	\[
		f(x)=x^m g(x),\qquad
		g(x)=\frac1{(m-1)!}\int_0^1(1-t)^{m-1}f^{(m)}(tx)\,\mathrm dt.
	\]
	积分区间固定且被积函数光滑，故 \(g\) 在 \(0\) 的一个邻域光滑。因此 \(\ker J\) 的每个茎都等于乘 \(x^m\) 的像，\cref{g2:prop:stalk-exactness}便给出所需短正合列与余核同构。余核在 \(0\) 处的茎为 \(\mathbb R^m\)，在其余点为零；常值层 \(\underline{\mathbb R^m}\) 的每点茎都为 \(\mathbb R^m\)，所以两者不同。这个例子把“局部无法除以 \(x^m\)”精确压缩为有限阶导数资料。
\end{solution}

\begin{exercise}\label{g2:ex:cech-cycle-refinement}
	设圆有一个 \(N\ge3\) 个开弧组成的覆盖，只有相邻两弧相交，所有非空交连通，没有三重交。对常值层 \(\underline R\)，按绕圆的方向给各相邻交定向，把一次上链写成 \((c_0,\ldots,c_{N-1})\)。证明一次 Čech 类由 \(\sum_i c_i\) 唯一确定。再设一个 \(2N\) 弧覆盖细化它，指标选择按模 \(N\) 记为 \(r(2i)=i\)、\(r(2i+1)=i+1\)；计算细化对上述和的影响，并说明反转圆的方向产生什么变化。
\end{exercise}
\begin{solution}
	零次上链为 \((b_i)\in R^N\)，一次微分在定向的相邻交上为 \((\delta b)_i=b_{i+1}-b_i\)，指标按模 \(N\) 取值。这仍是交错 Čech 约定，只是最后一条边采用从 \(N-1\) 指向 \(0\) 的方向，而不是指标递增方向。所有余边界的分量和为零。反之若 \(\sum_i c_i=0\)，取 \(b_0=0\)，并依次令 \(b_{i+1}=b_i+c_i\)；总和为零保证最后回到 \(b_N=b_0\)。于是 \(c=\delta b\)。故求和诱导 \(\check H^1\cong R\)，同时零次闭元为共同常数。
	在细化覆盖上，从 \(2i\) 到 \(2i+1\) 的边映到旧覆盖的 \(i\to i+1\) 边，故取值为 \(c_i\)；从 \(2i+1\) 到 \(2i+2\) 的边，两端旧指标同为 \(i+1\)，按重复指标约定取值为零。因此细化后的分量为 \((c_0,0,c_1,0,\ldots,c_{N-1},0)\)，总和不变，细化诱导此 \(R\) 上的恒等映射。若反转绕行方向，各边取负且次序反转，总和变号。这里识别一次群所用的方向是额外选择，而群及细化映射本身并未依赖该选择。
\end{solution}

\begin{exercise}\label{g2:ex:transverse-face-extension}
	在坐标轴的并集上给定 \(b(x,0)=x^2\)、\(b(0,y)=y^3\)，写出一个光滑延拓。再考虑直线 \(L=\{(x,0)\}\) 与抛物线 \(P=\{(x,x^2)\}\)，在其并集上给定 \(c(x,0)=0\)、\(c(x,x^2)=x\)。说明 \(c\) 连续且在每条曲线上光滑，却不能延拓为原点附近的 \(C^1\) 函数。解释这与\cref{g2:lem:simplex-face-extension}并不矛盾。
\end{exercise}
\begin{solution}
	第一组数据由 \(B(x,y)=x^2+y^3\) 延拓。两条坐标轴相交于原点，减去被重复计算的公共值后，正好得到\cref{g2:eq:corner-face-extension}的二维形式。

	第二组数据在两曲线的唯一交点都等于零。若并集中的点趋于原点，它的横坐标趋于零，两种公式的值都趋于零；在其他点附近只有一条曲线需要考虑，所以 \(c\) 连续。参数 \(x\mapsto(x,0)\) 与 \(x\mapsto(x,x^2)\) 都是光滑嵌入，沿它们的限制分别为 \(0\) 和 \(x\)，故逐曲线光滑。假设有 \(C^1\) 延拓 \(C\)，其 \(\partial_yC\) 在原点附近有界，于是
	\[
		x=C(x,x^2)-C(x,0)=\int_0^{x^2}\partial_yC(x,t)\,\mathrm{d}t=O(x^2),
	\]
	这在 \(x\to0\)、\(x\ne0\) 时不可能。单形的相邻面可以同时写成不同坐标超平面，法向线性独立；这里的两条曲线却相切，不能用同一坐标变换化为两条坐标轴。因此，逐片光滑的数据能否延拓还取决于这些片的相交方式。
\end{solution}

\begin{exercise}\label{g2:ex:close-circle-loops}
	把 \(S^1\) 看作复平面内的单位圆。设光滑闭路 \(\gamma,\eta:[0,1]\to S^1\) 满足 \(\gamma(0)=\gamma(1)=\eta(0)=\eta(1)=1\)，并且 \(\sup_t|\gamma(t)-\eta(t)|<2\)。构造保持端点的逐棱柱光滑同伦，并证明两闭路的绕数相同。再说明常数 \(2\) 不能放宽为允许等号。
\end{exercise}
\begin{solution}
	两个单位复数之间的线段经过原点，当且仅当它们互为相反数；此时两者距离恰为 \(2\)。所以题设使
	\[
		H(t,s)=\frac{(1-s)\gamma(t)+s\eta(t)}{|(1-s)\gamma(t)+s\eta(t)|}
	\]
	处处有定义。分母在紧正方形上有正下界，因而这个公式在正方形附近有光滑延拓，尤其在任意棱柱剖分的各单形上光滑；\(t=0,1\) 时恒为 \(1\)，故端点固定。取单位圆上的闭形式 \(\alpha=(x\,\mathrm{d}y-y\,\mathrm{d}x)/(2\pi)\)。给正方形取 \(\mathrm{d}s\wedge\mathrm{d}t\) 定向，剖分为两个有向三角形并用 Stokes 公式，内部边抵消，两个端点侧边为常值，得到
	\[
		\int_\eta\alpha-\int_\gamma\alpha=\int_{[0,1]^2}H^*(\mathrm{d}\alpha)=0.
	\]
	这两个积分就是绕数。若取 \(\gamma(t)=e^{2\pi\mathrm{i}t}\)、\(\eta(t)=1\)，最大距离为 \(2\)，而绕数分别为 \(1\) 与 \(0\)。它们不可能相对端点同伦：第一卷定理5.1.13证明，以 \(1\) 为基点的圆周环路的定端点同伦类由绕数完全分类。这个例子说明，欧氏空间中的线性插值必须留在回缩的定义域内；仅仅写出插值公式还没有得到目标流形内的同伦。
\end{solution}

\begin{exercise}\label{g2:ex:integration-not-strict-product}
	在 \(\mathbb R^2\) 上取 \(\omega=\mathrm dx\)、\(\eta=\mathrm dy\)，并令 \(\sigma:\Delta^2\to\mathbb R^2\) 是顶点为 \((0,0),(1,0),(0,1)\) 的标准仿射单形。计算 \(I(\omega\wedge\eta)(\sigma)\) 与 \((I\omega\smile I\eta)(\sigma)\)。再构造一个一次光滑奇异上链 \(T\)，使 \(I\omega\smile I\eta-I(\omega\wedge\eta)=\delta T\)，并说明这与正文的乘法定理相容。
\end{exercise}
\begin{solution}
	三角形有向面积为 \(1/2\)，故第一值为 \(1/2\)。前后面杯积给出 \(\int_{01}\mathrm dx\cdot\int_{12}\mathrm dy=1\cdot1=1\)。这说明积分不是严格的微分分次代数同态。

	令 \(a=I(x)\) 为零上链，即在点上取其 \(x\) 坐标，定义
	\[
		T=a\smile I(\mathrm dy)-I(x\,\mathrm dy).
	\]
	积分与微分相容，杯积满足 Leibniz 公式，且 \(\delta I(\mathrm dy)=0\)，所以
	\[
		\delta T=I(\mathrm dx)\smile I(\mathrm dy)-I(\mathrm dx\wedge\mathrm dy).
	\]
	具体地，对一条光滑道路 \(\tau\)，\(T(\tau)=x(\tau(0))(y(\tau(1))-y(\tau(0)))-\int_\tau x\,\mathrm dy\)。在边 \(01,12,02\) 上依次为 \(0,1/2,0\)，故 \((\delta T)(\sigma)=1/2\)，与前述差相符。两种上链不相等，却相差余边界，因而上同调乘法仍相同。
\end{solution}

\begin{exercise}\label{g2:ex:relative-cone-line}
	取 \(M=\mathbb R\)、\(V=\mathbb R\setminus\{0\}\)。在相对形式复形 \(A^*(M,V)\) 中，令 \(\eta\) 在负半轴为 \(0\)、在正半轴为 \(1\)，考虑一次圈 \((0,\eta)\)。证明其类非零，求它在从 \(-1\) 到 \(1\) 的有向相对区间类上的值，并用紧支形式表示它在 \(\varinjlim_KH^1(A^*(\mathbb R,\mathbb R\setminus K))\) 中的像。
\end{exercise}
\begin{solution}
	\(\eta\) 在 \(V\) 上局部常值，故 \(\mathrm D(0,\eta)=(0,-\mathrm d\eta)=0\)。若它是边界，则存在光滑函数 \(f\) 满足 \(\mathrm df=0\)、\(f|_V=\eta\)。前者迫使 \(f\) 在连通实轴上为一个常数，后者却要求两个不同常数，矛盾。

	把 \(\eta\) 的零上链任意延拓为实轴上的零上链 \(b\)。相对比较将 \((0,I\eta)\) 与 \((-\delta b,0)\) 识别，所以在区间上的值是 \(-b(1)+b(-1)=-1\)。负号来自本章相对复形的约定 \(\mathrm D(\omega,\eta)=(\mathrm d\omega,\omega|-\mathrm d\eta)\)，不能只凭两端值之差猜它的顺序。

	取光滑函数 \(\chi\)，在 \(( -\infty,-1]\) 为 \(0\)，在 \([1,\infty)\) 为 \(1\)。当紧集扩大到包含 \([-1,1]\) 时，\(\chi|=\eta\)，故减去 \(\mathrm D(\chi,0)\) 把原圈变为 \((-\mathrm d\chi,0)\)。其第一项紧支，且 \(\int_{\mathbb R}-\mathrm d\chi=-1\)，表示所求非零紧支类。
\end{solution}

\begin{exercise}\label{g2:ex:proper-pullback-compact-forms}
	用积分将 \(H_{c,\mathrm{dR}}^1(\mathbb R)\) 识别为 \(\mathbb R\)。计算固有光滑映射 \(f(x)=x^2\)、\(g(x)=-x\) 在该群上的拉回。再说明覆盖映射 \(p:\mathbb R\to S^1\)、\(p(t)=(\cos2\pi t,\sin2\pi t)\) 为什么不能以通常拉回来定义从全部紧支形式到紧支形式的映射。
\end{exercise}
\begin{solution}
	紧支一次形式写成 \(h(y)\,\mathrm dy\)。若积分为零，其原函数 \(F(y)=\int_{-\infty}^yh(s)\,\mathrm ds\) 在两端均为零，故紧支；任意实数又可由紧支函数的积分实现。因此积分确实给出所述识别。

	\(f^*(h\,\mathrm dy)=2x h(x^2)\,\mathrm dx\)，这是紧支奇函数系数的一次形式，其全轴积分为零，故 \(f^*=0\)。还可直接取 \(G(y)=-\int_y^\infty h(s)\,\mathrm ds\)，则 \(G(x^2)\) 在 \(|x|\) 足够大时为零，并满足 \(\mathrm d(G(x^2))=2x h(x^2)\,\mathrm dx\)。对 \(g\)，拉回为 \(-h(-x)\,\mathrm dx\)，其积分为 \(-\int h\)，故 \(g^*\) 是乘以 \(-1\)。

	圆周紧，所以其每个形式都紧支，但归一化角形式 \(\theta=(x\,\mathrm dy-y\,\mathrm dx)/(2\pi)\) 的拉回为 \(\mathrm dt\)，支集是整条实轴。\(p\) 的确不是固有映射：紧集 \(S^1\) 的逆像不紧。因此普通 de Rham 拉回存在，紧支拉回却不在所述复形之间定义。
\end{solution}

\begin{exercise}\label{g2:ex:integral-one-form-circle-map}
	设 \(M\) 为有限维、Hausdorff、第二可数的连通无边界光滑流形，\(\omega\in\Omega^1(M)\) 为闭实一形式。记 \(\vartheta\) 为 \(\mathbb R/\mathbb Z\) 上由 \(\mathrm dt\) 下降的一形式。证明：\(\omega\) 在每条分片光滑闭路上的积分均为整数，当且仅当存在光滑映射 \(f:M\to\mathbb R/\mathbb Z\)，使 \(\omega=f^*\vartheta\)。说明在固定一点的 \(f\) 值后，这个映射是否唯一。
\end{exercise}
\begin{solution}
	先设周期为整数。固定 \(p\in M\) 及 \(c\in\mathbb R/\mathbb Z\)，对从 \(p\) 到 \(x\) 的分片光滑道路 \(\gamma\)，定义
	\[
		f(x)=c+\left[\int_\gamma\omega\right]\pmod{\mathbb Z}.
	\]
	连通光滑流形可由坐标球连接任意两点，故这样的道路存在。两条道路的积分差为闭路 \(\gamma*\bar\gamma'\) 的积分，按假设为整数，因此 \(f\) 良定义。在小坐标球 \(U\) 上，Poincaré 引理给出 \(\omega|_U=\mathrm dh\)。取 \(x_0\in U\)，并把到 \(x\) 的道路分成固定的 \(p\) 到 \(x_0\) 一段及 \(U\) 内一段，便有 \(f(x)=f(x_0)+[h(x)-h(x_0)]\)。所以 \(f\) 局部具有光滑实值提升，并且 \(f^*\vartheta=\mathrm dh=\omega\)。

	反过来，若 \(\omega=f^*\vartheta\)，对一条闭路 \(\gamma\)，将 \(f\circ\gamma\) 提升到实线，记提升为 \(\widetilde\gamma\)。端点之差是整数，且
	\[
		\int_\gamma\omega=\int_{f\circ\gamma}\vartheta
		=\widetilde\gamma(1)-\widetilde\gamma(0)\in\mathbb Z.
	\]
	若 \(f_1,f_2\) 的拉回均为 \(\omega\)，则群值差 \(f_1-f_2\) 的每个局部实提升都具有零微分，故局部常值。\(M\) 连通使它整体为常值；固定 \(p\) 处的值后，差为零，因而唯一。这个题目表明整周期的一形式可能没有实值原函数，却有圆周值原函数。
\end{solution}

\begin{exercise}\label{g2:ex:dense-period-group}
	在 \(T^2=\mathbb R^2/\mathbb Z^2\) 上取 \(\omega=\mathrm dx+\sqrt2\,\mathrm dy\)。求其周期群 \(\operatorname{Per}_{\omega}(H_1(T^2;\mathbb Z))\)，证明它在 \(\mathbb R\) 中稠密，并证明任何非零实数倍 \(a\omega\) 都不具有整周期。
\end{exercise}
\begin{solution}
	两条坐标圆周生成一次整数同调，周期分别为 \(1,\sqrt2\)，所以周期群为 \(G=\mathbb Z+\sqrt2\mathbb Z\)。要证稠密，给定整数 \(N\geq1\)，考察 \(0,\sqrt2,\ldots,N\sqrt2\) 的小数部分。把 \([0,1)\) 分成 \(N\) 个等长区间，必有两个小数部分落在同一区间内。因此存在非零整数 \(q\)、整数 \(p\)，使 \(0<|q\sqrt2-p|<1/N\)；严格正性来自 \(\sqrt2\) 无理。适当变号，得到 \(G\) 中一个 \(0<\varepsilon<1/N\)。任给 \(r\in\mathbb R\)，取最接近 \(r/\varepsilon\) 的整数 \(j\)，则 \(|r-j\varepsilon|\leq\varepsilon/2\)。令 \(N\) 任意大，证明 \(G\) 稠密。

	若 \(a\omega\) 具有整周期，两条坐标圆周要求 \(a\in\mathbb Z\) 且 \(a\sqrt2\in\mathbb Z\)。当 \(a\ne0\) 时，两者之比会使 \(\sqrt2\) 为有理数，矛盾。周期群可以稠密，这并不与\cref{g2:prop:integral-periods}的格结论冲突：后者说的是具有整周期的上同调类在整个上同调向量空间中组成一个格，而不是说任一实上同调类的周期值总组成离散群。
\end{solution}

\begin{exercise}\label{g2:ex:projective-power-map-forms}
	对整数 \(d\geq1\)、\(m\geq1\)，定义
	\[
		F_d:\mathbb CP^m\longrightarrow\mathbb CP^m,
		\qquad [z_0:\cdots:z_m]\longmapsto[z_0^d:\cdots:z_m^d].
	\]
	证明此映射良定义且光滑。若 \(\omega\) 是\eqref{g2:eq:projective-de-rham-ring}的归一化闭二形式，求 \([F_d^*\omega]\) 及 \(F_d\) 在最高次上同调上的作用，从而求其映射度。
\end{exercise}
\begin{solution}
	同乘一个非零复数 \(\lambda\) 后，各坐标的 \(d\) 次幂同乘 \(\lambda^d\)，故射影点不变；至少一个输入坐标非零，输出也不会全为零。在第 \(j\) 个输入坐标非零的仿射图上，输出的同一个坐标非零，局部表达为各仿射坐标逐个取 \(d\) 次幂，所以光滑。

	坐标直线 \(\mathbb CP^1\) 在此映射下不变，其仿射表达为 \(z\mapsto z^d\)。值 \(1\) 有恰好 \(d\) 个原像；每个原像处复导数 \(dz^{d-1}\) 非零，相应实 Jacobian 的行列式为 \(|dz^{d-1}|^2>0\)。由\cref{g2:thm:local-degree-sum}，这个球面映射的度为 \(d\)。因此 \(F_d^*\omega\) 在正向坐标直线上的积分为 \(d\)，而 \(\omega\) 的积分为 \(1\)。二次上同调是一维，故
	\[
		[F_d^*\omega]=d[\omega].
	\]
	拉回保持楔积，所以 \(F_d^*[\omega^j]=d^j[\omega^j]\)。特别地，\([\omega^m]\ne0\) 生成顶次实上同调，\(F_d^*\) 在该处为乘 \(d^m\)。用与顶次基本类的自然评价，顶次上同调上的乘数就是映射度，于是 \(\deg F_d=d^m\)。这里仅用了闭形式的类；一般没有 \(F_d^*\omega=d\omega\) 的逐点等式。
\end{solution}

\begin{exercise}\label{g2:ex:curvature-obstructs-complex}
	在 \(\mathbb R^2\) 的平凡实线丛上取联络 \(\nabla f=\mathrm df+x f\,\mathrm dy\)。计算 \(\mathrm d_\nabla^2f\)，判断 \(\operatorname{im}(\nabla:\Omega^0\to\Omega^1)\) 是否包含在 \(\ker(\mathrm d_\nabla:\Omega^1\to\Omega^2)\) 中，并求任一非空连通开集上的全部平行截面。
\end{exercise}
\begin{solution}
	联络矩阵为 \(A=x\,\mathrm dy\)，故 \(F=\mathrm dA+A\wedge A=\mathrm dx\wedge\mathrm dy\)，从而
	\[
		\mathrm d_\nabla^2f=f\,\mathrm dx\wedge\mathrm dy.
	\]
	例如 \(\nabla1=x\,\mathrm dy\)，其外协变微分为 \(\mathrm dx\wedge\mathrm dy\ne0\)，所以像不包含于核，不能用这两个空间的商定义上同调。若 \(\nabla f=0\)，再作用一次 \(\mathrm d_\nabla\) 得 \(f\,\mathrm dx\wedge\mathrm dy=0\)；面积形式处处非零，因此 \(f=0\)。这在每个非空开集上都成立，说明非平坦联络甚至可能没有局部非零平行截面，不能把它的平行截面当成与纤维同秩的局部系统。
\end{solution}

\begin{exercise}\label{g2:ex:flat-torus-contracting-homotopy}
	在 \(T^2=\mathbb R^2/\mathbb Z^2\) 的平凡实线丛上取
	\[
		\mathrm d_\nabla=\mathrm d+(a\,\mathrm dx+b\,\mathrm dy)\wedge,
		\qquad (a,b)\in\mathbb R^2\setminus\{(0,0)\}.
	\]
	证明它给出一个平坦联络，并构造该形式复形的收缩同伦，从而证明所有次数的上同调均为零。
\end{exercise}
\begin{solution}
	一形式 \(A=a\,\mathrm dx+b\,\mathrm dy\) 闭，且在线丛中 \(A\wedge A=0\)，故曲率为零。交换坐标后可设 \(a\ne0\)。在光滑双周期函数上令 \(D_x=\partial_x+a\)、\(D_y=\partial_y+b\)。对给定 \(f\)，方程 \(D_xu=f\) 的周期解唯一，且为
	\[
		u(x,y)=e^{-ax}\left(\frac{\int_0^1e^{at}f(t,y)\,\mathrm dt}{e^a-1}
		+\int_0^x e^{at}f(t,y)\,\mathrm dt\right).
	\]
	由于 \(a\) 为非零实数，分母不为零。公式给出 \(u(1,y)=u(0,y)\)，微分方程及 \(f\) 的周期性进一步给出各阶导数的端点一致性；对 \(y\) 的光滑性和周期性可直接在公式中检查。齐次解 \(ce^{-ax}\) 只有在 \(c=0\) 时周期，所以解唯一。记此逆算子为 \(B=D_x^{-1}\)，并逐分量作用于形式。它与 \(D_x,D_y\) 交换：对 \(D_x\) 由互逆关系成立，对 \(D_y\) 则由公式中参数微分与积分交换成立。

	在形式上有 \(\mathrm d_\nabla=\mathrm dx\wedge D_x+\mathrm dy\wedge D_y\)。定义降一次算子 \(K=B\iota_{\partial_x}\)。收缩算子满足
	\[
		\iota_{\partial_x}(\mathrm dx\wedge\eta)+\mathrm dx\wedge\iota_{\partial_x}\eta=\eta,
		\qquad
		\iota_{\partial_x}(\mathrm dy\wedge\eta)+\mathrm dy\wedge\iota_{\partial_x}\eta=0.
	\]
	结合 \(BD_x=1\) 及所述交换性，得到 \(\mathrm d_\nabla K+K\mathrm d_\nabla=1\)。因此每个闭形式 \(\eta\) 都满足 \(\eta=\mathrm d_\nabla K\eta\)，全部上同调为零。这个平坦丛的底层线丛虽平凡，指定联络沿两条坐标圆周的运输却为 \(e^{-a}\)、\(e^{-b}\)；至少一个不是 \(1\)，所以其微分形式上同调可以与普通环面上同调完全不同。
\end{solution}
